\documentclass[12pt, a4paper, reqno, oneside]{book}


\usepackage[left=2.5cm, right=2cm, top=2cm, bottom=2cm]{geometry}
\usepackage{amsmath,amsxtra,amssymb,latexsym, amscd,amsthm}
\usepackage{microtype}
\usepackage{indentfirst}
\usepackage{color}
\usepackage{empheq}
\usepackage[USenglish]{babel}
\usepackage{amsmath}
\usepackage[mathscr]{eucal}
\usepackage{amsfonts}
\usepackage{graphicx}
\usepackage{float}
\usepackage{listings}
\usepackage[dvipsnames]{xcolor}
\usepackage{verbatim}
\usepackage{tikz}
\usetikzlibrary{calc}
\usepackage{tkz-graph}
\usetikzlibrary {positioning}
\usepackage{float}
\usepackage{array}
\usepackage{multirow}
\usepackage{algorithm}
\usepackage{algpseudocode}
\usepackage{subfigure}
\usepackage{pgfplots}
\usepackage{titlesec}
\usepackage{fancyhdr}
\usepackage{rotating}
\usepackage{longtable}
\usepackage{supertabular}
\fancyhf{}
% \fancyhead[LO,RE]{\small{\leftmark}}
\fancyfoot[C]{ \thepage}
\usepackage{ulem}
\usepackage{pdfpages}
\usepackage{niceframe}
\usepackage[square,numbers]{natbib}
\usepackage{subcaption,booktabs}
\usepackage{pgfplots}
\pgfplotsset{compat=1.18}
\usepackage{chngcntr}
\usepackage{tabularray}
\usepackage[labelfont=bf, textfont=it, justification=centering]{caption}
\captionsetup{labelfont={color=blue, bf}, textfont={color=black, it}}

\hyphenation{op-tical net-works semi-conduc-tor}
%\input setbmp
\usepackage[hyphens,spaces,obeyspaces]{url}
\usepackage[colorlinks,bookmarksopen,bookmarksnumbered,citecolor=blue, colorlinks=true,linkcolor=blue,urlcolor=blue]{hyperref}
% \usepackage{url}

% \usepackage{enumitem}
%\usepackage{titledot}
%\titlename{Chương}
%+ formating
%\setlength{\oddsidemargin}{0in}
%\setlength{\textwidth}{6.in}
%\setlength{\topmargin}{-0.5in}
%\setlength{\textheight}{9.25in}
\renewcommand{\baselinestretch}{1.25}

%Chương không số
\newcommand \chap[1]{%
  \chapter*{#1}%
  \addcontentsline{toc}{chapter}{#1}
  \markboth{#1}{}
}
\newcommand \sect[1]{%
  \section*{#1}%
  \addcontentsline{toc}{section}{#1}}
  
%+ Dinh nghia mot so ky hieu toan hoc.
\def\B{\mathscr B}
\def\one{\mathbf 1}
\def\R{{\mathbb R}}
\def\LL{{\mathbb L}}
\newcommand\abs[1]{|#1|}
\newcommand\set[1]{\{#1\}}
\newcommand\norm[1]{||#1||}
\DeclareMathOperator{\vrai}{vrai}
\DeclareMathOperator{\const}{const}
\newcommand{\eoproof}{\hfill $\square$}
\newcommand{\spro}[1]{\left<#1\right>}
\newcommand{\argmin}{\arg\min}

\counterwithin{algorithm}{chapter}
\renewcommand{\thealgorithm}{\thechapter.\arabic{algorithm}}
\renewcommand{\theequation}{\arabic{chapter}.\arabic{equation}} %\arabic{section}.
\newtheorem{lemma}{\bf Bổ đề}[chapter]
\newtheorem{theorem}{\bf Định lý}[chapter]
\newtheorem{proposition}{\bf Mệnh đề}[chapter]
\newtheorem{corollary}{\bf Hệ quả}[chapter]
\newtheorem{definition}{\bf Định nghĩa}[chapter]
\newtheorem{remark}{\bf Chú ý}[section]
% \newtheorem{problem}{Bài toán}
% \renewcommand*{\listalgorithmcfname}{List of Algorithms}
% \renewcommand*{\algorithmcfname}{Algorithm}
% \renewcommand*{\algorithmautorefname}{thuật toán}

%mức tiêu đề
\titleclass{\subsubsubsection}{straight}[\subsection]

\newcounter{subsubsubsection}[subsubsection]
\renewcommand\thesubsubsubsection{\thesubsubsection.\arabic{subsubsubsection}}
\renewcommand\theparagraph{\thesubsubsubsection.\arabic{paragraph}} % optional; useful if paragraphs are to be numbered

\titleformat{\subsubsubsection}
  {\normalfont\normalsize\bfseries}{\thesubsubsubsection}{1em}{}
\titlespacing*{\subsubsubsection}
{0pt}{3.25ex plus 1ex minus .2ex}{1.5ex plus .2ex}

\makeatletter
\renewcommand\paragraph{\@startsection{paragraph}{5}{\z@}%
  {3.25ex \@plus1ex \@minus.2ex}%
  {-1em}%
  {\normalfont\normalsize\bfseries}}
\renewcommand\subparagraph{\@startsection{subparagraph}{6}{\parindent}%
  {3.25ex \@plus1ex \@minus .2ex}%
  {-1em}%
  {\normalfont\normalsize\bfseries}}
\def\toclevel@subsubsubsection{4}
\def\toclevel@paragraph{5}
\def\toclevel@paragraph{6}
\def\l@subsubsubsection{\@dottedtocline{4}{7em}{4em}}
\def\l@paragraph{\@dottedtocline{5}{10em}{5em}}
\def\l@subparagraph{\@dottedtocline{6}{14em}{6em}}
\makeatother

\setcounter{secnumdepth}{4}
\setcounter{tocdepth}{4}


%thư mục hình
\graphicspath{{./figs/}}

\definecolor{mygreen}{rgb}{0,0.6,0}
\definecolor{mygray}{rgb}{0.5,0.5,0.5}
\definecolor{mymauve}{rgb}{0.58,0,0.82}

\lstdefinestyle{customc}{
  belowcaptionskip=1\baselineskip,
  breaklines=true,
  frame=single,
  xleftmargin=\parindent,
  language=C,
  showstringspaces=false,
  basicstyle=\footnotesize\ttfamily,
  keywordstyle=\bfseries\color{green!40!black},
  commentstyle=\itshape\color{purple!40!black},
  identifierstyle=\color{blue},
  stringstyle=\color{orange},
  numbers=left,                   % where to put the line-numbers
  numberstyle=\tiny\color{gray},  % the style that is used for the line-numbers
  stepnumber=1, 
}

\lstdefinestyle{customasm}{
  belowcaptionskip=1\baselineskip,
  frame=single,
  xleftmargin=\parindent,
  language=[x86masm]Assembler,
  basicstyle=\footnotesize\ttfamily,
  commentstyle=\itshape\color{purple!40!black},
  numbers=left,                   % where to put the line-numbers
  numberstyle=\tiny\color{gray},  % the style that is used for the line-numbers
  stepnumber=1, 
}

\lstset{escapechar=@,style=customc}

\definecolor {processblue}{cmyk}{0.96,0,0,0}

\newcolumntype{C}[1]{>{\centering\arraybackslash}p{#1}}

\definecolor{eclipseStrings}{RGB}{42,0.0,255}

\definecolor{eclipseKeywords}{RGB}{127,0,85}

\colorlet{numb}{magenta!60!black}



\lstdefinelanguage{json}{
    basicstyle=\normalfont\ttfamily,
    commentstyle=\color{eclipseStrings}, % style of comment
    stringstyle=\color{eclipseKeywords}, % style of strings
    numbers=left,
    numberstyle=\scriptsize,
    stepnumber=1,
    numbersep=8pt,
    showstringspaces=false,
    breaklines=true,
    frame=lines,
    string=[s]{"}{"},
    comment=[l]{:\ "},
    morecomment=[l]{:"},
    literate=
        *{0}{{{\color{numb}0}}}{1}
         {1}{{{\color{numb}1}}}{1}
         {2}{{{\color{numb}2}}}{1}
         {3}{{{\color{numb}3}}}{1}
         {4}{{{\color{numb}4}}}{1}
         {5}{{{\color{numb}5}}}{1}
         {6}{{{\color{numb}6}}}{1}
         {7}{{{\color{numb}7}}}{1}
         {8}{{{\color{numb}8}}}{1}
         {9}{{{\color{numb}9}}}{1}
}

\begin{document}
\hypersetup{pageanchor=false}

\begin{titlepage}
\begin{center}
	\begin{tikzpicture}[overlay,remember picture]
	    \draw [line width=3pt,rounded corners=0pt,
	        ]
	        ($ (current page.north west) + (20mm,-20mm) $)
	        rectangle
	        ($ (current page.south east) + (-15mm,20mm) $);
	    \draw [line width=1pt,rounded corners=0pt]
	        ($ (current page.north west) + (21.5mm,-21.5mm) $)
	        rectangle
	        ($ (current page.south east) + (-16.5mm,21.5mm) $);
	\end{tikzpicture}
	\\[1mm]
    \centerline{\bf VIETNAM NATIONAL UNIVERSITY HANOI}
    \centerline{\bf UNIVERSITY OF ENGINEERING AND TECHNOLOGY}
    
    \vspace*{1.8cm}
    \large
        \begin{figure}[H]
        \centering
        
\includegraphics[width=0.23\linewidth]{figs/uet.jpg}
        \end{figure}
    \vspace*{1cm}
    \centerline{\bf TRAN VAN SON}
    \vspace*{2.5cm}
    \centerline{\bf  A SECURE FRAMEWORK FOR ENFORCING}
    \centerline{\bf   ZERO TRUST POLICIES IN DISAGGREGATED NETWORKS }
     
    % \centerline{\bf A Secure Framework for Enforcing Zero Trust Policies in NFV Environments}
    \centerline{\bf}
    \vspace*{2.5cm}
    \centerline{\bf BACHELOR THESIS}
    \vspace*{0.2cm}
    \centerline{\bf Major: Computer Science}
    
    \vfill
    \centerline{\bf Hanoi - 2025}
    \vspace{13pt}
\end{center}
\end{titlepage}

\newpage

\begin{titlepage}
\begin{center}
	\begin{tikzpicture}[overlay,remember picture]
	    \draw [line width=3pt,rounded corners=0pt,
	        ]
	        ($ (current page.north west) + (20mm,-20mm) $)
	        rectangle
	        ($ (current page.south east) + (-15mm,20mm) $);
	    \draw [line width=1pt,rounded corners=0pt]
	        ($ (current page.north west) + (21.5mm,-21.5mm) $)
	        rectangle
	        ($ (current page.south east) + (-16.5mm,21.5mm) $);
	\end{tikzpicture}
	\\[1mm]
    \centerline{\bf VIETNAM NATIONAL UNIVERSITY HANOI}
    \centerline{\bf UNIVERSITY OF ENGINEERING AND TECHNOLOGY}
    \vspace*{1.8cm}
    \large
        \begin{figure}[H]
        \centering
        
\includegraphics[width=0.23\linewidth]{figs/uet.jpg}
        \end{figure}
    \vspace*{0.9cm}
    \large
    \centerline{\bf TRAN VAN SON}
    \vspace*{2.2cm}
    \centerline{\bf  A SECURE FRAMEWORK FOR ENFORCING}
    \centerline{\bf   ZERO TRUST POLICIES IN DISAGGREGATED NETWORKS }
    \vspace*{1.7cm}
    \centerline{\bf Major: Computer Science}
    
    \vspace*{1.6cm}
    \centerline{\bf Supervisor: }
    \centerline{\bf Assoc. Prof. Nguyen Ngoc Hoa }
    
    \vfill
    \centerline{\bf Hanoi - 2025}
    \vspace{13pt}
\end{center}
\end{titlepage}



\newpage
\pagenumbering{gobble}
\chapter*{Acknowledgements}
\vspace{20pt}
% Fill the acknowledgements

I would like to express my sincere gratitude to my supervisor, Associate Professor Nguyen Ngoc Hoa, for his invaluable guidance and encouragement during this research. My professor's expertise and feedback were instrumental throughout the writing process, which made the completion of this thesis possible.

I acknowledge that my thesis still has potential for enhancement. I welcome any feedback that could help improve the quality of this thesis in the future.

Sincere gratitude!


\newpage
\chapter*{Authorship}
\vspace{20pt}
% Fill the authorship

I declare that the thesis is my original work, under the supervision of Associate Professor Nguyen Ngoc Hoa. This work has not been submitted for any other degree or professional qualification.

All references and related works used in this thesis have been appropriately acknowledged and cited. The experimental results and performance metrics presented herein are derived from rigorous testing and validation of the proposed framework within a controlled environment.

I affirm that this declaration is made with full integrity and in accordance with the institution's policies.

\begin{flushright}
{\it Author}
\hskip 1.5cm\quad


\vskip 2cm
{\bf Tran Van Son }  \quad\ 
\end{flushright}
\newpage



\newpage
% \pagenumbering{Roman}
\chapter*{Supervisor's Approval}
\vspace{20pt}
I hereby approve that the thesis in its current form is ready for committee examination as a requirement for the Bachelor of Computer Science degree at the VNU University of Engineering and Technology.

All references from relevant studies are explicitly sourced from the list of references in the thesis. The thesis does not copy documents and research works of others without specifying the references.


\begin{flushright}
{\it Signature}
\hskip 1.5cm\quad

\vskip 2cm
{\bf Nguyen Ngoc Hoa }  \quad\ 
\end{flushright}


\newpage



\centerline{\bf A SECURE FRAMEWORK FOR ENFORCING }
\centerline{\bf ZERO TRUST POLICIES IN DISAGGREGATED NETWORKS }
\vspace{10pt}
\centerline{Tran Van Son}

\vspace{2cm}
{\bf Abstract}
\vspace{5pt}

The telecommunications and data center industries are rapidly shifting to disaggregated network architectures, a paradigm that invalidates traditional perimeter-based security models and necessitates a move to a ZTA. However, the practical enforcement of Zero Trust policies is hindered by architectural and protocol incompatibilities, particularly between legacy orchestration systems like ONAP, which use NETCONF, and modern, cloud-native data planes like SONiC, which use gNMI.

This thesis addresses this critical security and integration gap by presenting the design, implementation, and evaluation of a \textbf{Secure Framework for Enforcing Zero Trust Policies in Disaggregated Networks}. The core of this work is a purpose-built \textbf{PEP} that functions as a secure, identity-aware intermediary. The framework integrates a granular \textbf{RBAC engine} to perform command-level authorization for every management request, ensuring that even authenticated orchestration systems operate under the principle of least privilege.

To provide defense-in-depth, the framework is architected to operate within a secure, identity-based overlay network (OpenZiti), which provides transport-layer security and architectural isolation by making management endpoints \textit{dark} to untrusted networks. The framework's internal logic performs high-fidelity translation between the NETCONF and gNMI paradigms, enabling consistent policy enforcement across heterogeneous devices without requiring modifications to existing systems.

The framework is  evaluated using a novel, hybrid methodology that combines a \textbf{Comparative Architectural Analysis} against the STRIDE threat model with an \textbf{Experimental Security Validation}. The results demonstrate that the proposed architecture provides a more complete security posture than existing single-layer security mechanisms and that its implemented controls function correctly. This work provides a foundational and practically validated blueprint for securing the next generation of automated, disaggregated networks.


% \centerline{\bf A SECURE FRAMEWORK FOR ENFORCING }
% \centerline{\bf ZERO TRUST POLICIES IN DISAGGREGATED NETWORKS }
% \vspace{10pt}


% \vspace{2cm}
% {\bf Abstract}
% \vspace{5pt}

% Fill the abstract



\newpage
%Table of Contents
\tableofcontents

\newpage
\listoffigures
 

\newpage
\listoftables
\begingroup
% \listofalgorithms
\let\clearpage\relax
\endgroup

\newpage
\chapter*{List of Abbreviations}
\vspace{20pt}

\begin{table}[H]
\centering
\begin{tabular}{|m{3cm}|m{8cm}|}
% Fill the abbreviations
\hline
ONAP & Open Network Automation Platform \\ \hline

SONiC  &  Software for Open Networking in the Cloud \\ \hline

gPC & Generalized Polynomial Chaos\\ \hline

VPN & Virtual Private Network \\ \hline
ZTA & Zero Trust Architecture \\ \hline

ZTNA & Zero Trust Network Access\\ \hline
RBAC &  Role-Based Access Control \\ \hline

NETCONF & Network Configuration Protocol \\ \hline

RESTCONF & Representational State Transfer Configuration Protocol \\ \hline
RPC & Remote Procedure Calls \\ \hline




mTLS & Mutual Transport Layer Security \\ \hline

% SNMP & Simple Network Management Protocol \\ \hline

gNMI & gRPC Network Management Interface \\ \hline

PEP & Policy Enforcement Point \\ \hline

SDN & Software-Defined Networking \\ \hline


NOS  & Network Operating System \\ \hline

YANG  & Yet Another Next Generation \\ \hline
\end{tabular}
\end{table}


\newpage
\pagestyle{fancy}
\pagenumbering{arabic}

\chap{Introduction}
\label{chap:introduction}

\sect{Motivation}
The telecommunications and data center industries are undergoing a fundamental transformation, shifting from monolithic, proprietary hardware to disaggregated architectures based on open-source software and white-box hardware \cite{borraccini2023disaggregated}. Two prominent open-source projects lead this evolution: the ONAP for service orchestration and the SONiC for the data plane. ONAP provides a policy-driven framework for automating the lifecycle of network services \cite{onap_architecture_2024, onap_evolution_2025}, while SONiC offers a flexible, high-performance Network Operating System (NOS) for modern data center appliances.

However, this disaggregated and dynamic environment dismantles traditional perimeter-based security models. With virtualized NOS instantiated and terminated dynamically, location-based trust becomes unreliable \cite{fernando2025addressing}. This necessitates a move towards a , a security model that operates on the principle of \textit{never trust, always verify} for every access request, regardless of its origin \cite{rose_zero_2020}. Implementing Zero Trust Architecture (ZTA) requires a robust mechanism to translate high-level security policies into granular, real-time enforcement actions at the data plane \cite{fernando2025addressing}.


The challenge of enforcing Zero Trust policies in these environments is compounded by a technical barrier: ONAP controllers communicate using NETCONF (RPC-based, XML-encoded) \cite{rfc6241_netconf}, while modern cloud-native systems like SONiC have adopted gNMI (stream-oriented, Protocol Buffers-encoded)  \cite{openconfig_gnmi_2018}. This protocol incompatibility prevents direct policy enforcement and creates a gap where security controls cannot be applied uniformly across the management plane. Without a mechanism to bridge this gap, Zero Trust policies cannot be consistently enforced across disaggregated network infrastructures.

This thesis addresses these challenges by introducing the Secure Framework, a purpose-built PEP designed specifically for enforcing Zero Trust policies in disaggregated networks. The framework serves as a secure intermediary that:
\begin{itemize}
    \item Enforces granular, command-level authorization through an integrated RBAC engine, ensuring that every management request-regardless of source-is evaluated against fine-grained policies before execution.
    
    \item Operates within a defense-in-depth architecture, leveraging identity-based secure transport mechanisms to ensure that policy enforcement occurs over authenticated, encrypted channels.
    
\item Bridges protocol heterogeneity by translating between management paradigms (NETCONF and gNMI), enabling consistent policy enforcement across diverse network elements without requiring changes to existing orchestration systems.
\end{itemize}


\sect{Research Questions}

\label{sec:research_challenges}
The development and evaluation of this Zero Trust policy enforcement framework address several key research questions as follows:

\begin{itemize}
    \item \textbf{RQ1: Fine-Grained Policy Enforcement in Dynamic Environments:} How can we enforce detailed Zero Trust policies in automated network systems where devices are constantly being added and removed?
    
    Traditional systems give broad permissions to authenticated controllers. This research designs a mechanism that enforces the principle of least privilege, giving only the minimum necessary permissions, by checking every management request against detailed policies. This prevents misuse by compromised or malfunctioning controllers, even when they are authenticated.
    
    \item \textbf{RQ2: Cross-Protocol Policy Consistency:} How can Zero Trust policies be consistently enforced when management operations use NETCONF while data plane operations use gNMI? 
    
    In this disaggregated architecture, the policy enforcement mechanism must translate and validate operations across these two protocols without compromising security guarantees. The challenge is designing an enforcement point that maintains semantic correctness when translating between NETCONF (management plane) and gNMI (data plane) while ensuring consistent policy application regardless of protocol differences.
    \item \textbf{RQ3: Architectural Isolation for Management Endpoints:} How can we isolate management endpoints to eliminate reconnaissance attacks while keeping them accessible to authorized systems? 
    
    This requires moving beyond traditional perimeter security to an identity-based model that enforces "authenticate before connect" workflows. Policy enforcement must occur over cryptographically secured channels, ensuring that endpoints are only visible to authorized entities.
    
    \item \textbf{RQ4: Evaluation of Defense-in-Depth Security Architecture:} How can we rigorously evaluate a multi-layered Zero Trust framework that combines isolated management endpoints with application-layer policy enforcement? 
    
    This requires implementing identity-based transport security that makes endpoints invisible to untrusted networks, and developing a comprehensive evaluation methodology. The evaluation combines comparative threat modeling (using STRIDE analysis) ~\cite{Khan2017stride, ruffy2016stride}  with empirical testing of integrated security controls across all architectural layers.
\end{itemize}

\sect{Thesis Objective and Contents}
\label{sec:objective}
The primary objective of this thesis is to design, implement, and comprehensively evaluate a \textbf{Zero Trust policy enforcement framework} for disaggregated networks. Specifically, this work demonstrates how 
\textbf{fine-grained, command-level access control} can be enforced on  network orchestration systems through a purpose-built PEP that operates over secure, identity-verified channels.


To achieve this objective, the thesis will cover the following contents:
\begin{itemize}
    \item A detailed analysis of \textbf{Zero Trust requirements} and \textbf{Policy enforcement challenges} in modern disaggregated networks, including the architectural and protocol barriers that prevent uniform security control.
    \item The design of the multi-layered security architecture, detailing the role of the secure overlay for transport security and the design of the Secure Framework for application-layer security, including its core translation logic and integrated RBAC engine.
    \item The development of a hybrid evaluation methodology, including a comparative architectural analysis against the STRIDE threat model and a set of experimental tests to validate the functional correctness of the security controls.
    \item A thorough discussion of the evaluation results, analyzing the framework's effectiveness and its implications for the future of secure network automation.
\end{itemize}

\sect{Highlight of Contributions}

This thesis introduces a novel framework for enforcing Zero Trust policies in disaggregated networks. The primary contributions are:
\begin{enumerate}
    \item \textbf{A Policy Enforcement Framework for Zero Trust in 
    Disaggregated Networks:} Design and implement the \textbf{Secure 
    Framework}, a purpose-built PEP that 
    enforces fine-grained, command-level access control on network 
    orchestration systems. The framework features:
    \begin{itemize}
        \item A \textbf{RBAC engine} that 
        evaluates every management request against granular policies, 
        enforcing the Zero Trust principle of least privilege even 
        for authenticated controllers.
        
        \item A \textbf{protocol-agnostic enforcement mechanism} that 
        operates consistently across heterogeneous management protocols 
        (NETCONF and gNMI), enabling uniform policy enforcement without 
        requiring modifications to existing orchestration systems.
        
        \item \textbf{Defense-in-depth integration} with identity-based 
        secure transport mechanisms, ensuring that policy enforcement 
        occurs over cryptographically authenticated channels where 
        management endpoints are architecturally isolated from 
        untrusted networks.
    \end{itemize}

    \item \textbf{Architectural Patterns for Zero Trust Implementation:} 
    Demonstrate a practical architectural pattern for implementing Zero 
    Trust in automated network environments, showing how application-layer 
    policy enforcement can be combined with transport-layer security 
    controls to achieve \textit{never trust, always verify} principles in 
    real-world disaggregated systems.

    \item \textbf{A Hybrid Evaluation Framework for Policy-Centric 
    Security:} Develop and apply a novel evaluation methodology 
    specifically designed for assessing policy enforcement mechanisms. 
    This framework combines:
    \begin{itemize}
        \item \textbf{Comparative threat modeling} using the STRIDE 
        framework to demonstrate architectural superiority over 
        single-layer security approaches.
        
        \item \textbf{Empirical validation} of policy enforcement 
        through systematic security testing, proving the functional 
        correctness of implemented access controls.
    \end{itemize}
\end{enumerate}

\sect{Thesis Outline}
\label{sec:thesis_outline}

This thesis comprises four primary chapters structured as follows.

\begin{itemize}
    \item Chapter \ref{chap:background} introduces the background information on the foundational technologies and concepts driving this research. 

    \item Chapter \ref{chap:related} outlines the related works in network security.

    \item Chapter \ref{chap:method} describes our proposed method in detail. It presents the Zero Trust architecture, explaining the integration of the OpenZiti overlay for secure transport. The chapter then provides an in-depth look at the design of the Secure Framework, covering its core translation logic and the implementation of its integrated security module.

    \item Chapter \ref{chap:experiments} presents the comprehensive evaluation of the proposed framework. It details the hybrid evaluation methodology, including the comparative architectural analysis against the STRIDE threat model and the experimental setup for security validation. 
\end{itemize}

Finally, the conclusion will summarize key contributions of our work and suggest potential directions for future research.


\chapter{Background}
\label{chap:background}

This chapter provides the foundational knowledge required to understand the core technologies and concepts underpinning this thesis. We begin by examining the industry-wide shift towards network disaggregation and the key open-source platforms driving this evolution: the ONAP and the SONiC. Subsequently, we delve into the technical specifics of the network management protocols central to this work, NETCONF and gNMI, highlighting their operational differences. Finally, we establish the security principles that motivate the thesis, focusing on the ZTA and ZTNA and its enabling technologies, namely mTLS and RBAC.

\section{Evolution of Network Architectures}
% \label{sec:overview}

The landscape of network infrastructure has undergone a paradigm shift over the past decade, moving away from vertically integrated, proprietary systems towards open, flexible, and disaggregated models. Historically, network operators were dependent on a small number of vendors who provided monolithic devices where the hardware, operating system, and control software were tightly coupled. This model stifled innovation, increased operational costs, and resulted in significant vendor lock-in. The disaggregation movement breaks this vertical stack, enabling network operators to independently select best-of-breed hardware and software components, fostering a more competitive and innovative ecosystem \cite{feamster2014case}. This evolution is primarily driven by two foundational technologies: SDN and Network Operation System (NOS).

\subsection{From Monolithic to Disaggregated Systems on Whitebox Hardware}

The core enabler of network disaggregation is the advent of "whitebox" hardware. These are generic, commodity switches and servers built with merchant silicon from manufacturers like Broadcom or Intel, rather than proprietary Application-Specific Integrated Circuits (ASICs). This hardware provides a standardized, open platform upon which various network operating systems and applications can be deployed. The decoupling of hardware from software allows for unprecedented flexibility, enabling cloud providers and telecommunication operators to customize their network stack to meet specific performance and functional requirements, often at a fraction of the cost of traditional systems \cite{Benson2010}. This disaggregation, however, places a greater burden on software to provide the control, management, and resilience previously handled by the monolithic vendor systems.

\subsection{Software-Defined Networking}
\label{sec:sdn}

SDN represents a paradigm shift in network architecture, founded on the principle of decoupling the network's control plane from its data plane. In traditional networks, the control logic is distributed across individual hardware devices, making global policy enforcement and dynamic management complex. SDN centralizes this control logic into a software-based entity known as the SDN Controller, which provides a programmable, holistic view of the entire network infrastructure \cite{paliwal2018controllers}.



The power of this approach is its ability to manage heterogeneous, multi-domain networks under a unified control system, as illustrated in Figure~\ref{fig:sdn_architecture}. This modern architecture extends beyond the data center to encompass the entire end-to-end service delivery path.

\begin{figure}[h!]
    \centering
    \includegraphics[width=\textwidth]{sdn_architecture.png} % Ensure this filename matches your uploaded image
    \caption{The SDN Architecture.}
    \label{fig:sdn_architecture}
\end{figure}

The architecture depicted in Figure~\ref{fig:sdn_architecture} is composed of several distinct, SDN-based domains:

\begin{itemize}
    \item \textbf{SDN-based Edge Network:} This is the access layer of the network, responsible for connecting end-user devices, sensors, and IoT hardware. The forwarding elements in this domain (e.g., wireless access points, edge switches) are programmable and receive their forwarding instructions from the central control plane.

    \item \textbf{SDN-based Core Network:} This domain serves as the high-speed backbone of the network. It interconnects the various edge networks and provides transport to the cloud services. The core routers and switches in this layer are critical for high-throughput packet forwarding and are managed by the SDN controller to optimize traffic paths and enforce network-wide policies. The secure, high-performance core router developed in this thesis is designed to operate within this domain.

    \item \textbf{SDN-based Cloud Network:} This layer hosts the applications, services, and aggregation servers that process the data generated by the edge devices. The connectivity and service chaining within this cloud environment are also managed via SDN principles, allowing for dynamic resource allocation and traffic steering.

    \item \textbf{SDN Control Plane:} This is the centralized intelligence of the architecture. As shown by the control-plane communication links (dashed lines), the controller has a global view and manages the behavior of both the Edge and Core networks via standardized Application Programming Interfaces (APIs), such as OpenFlow or, more recently, P4, which allows for programmable data plane behavior \cite{hauser2023survey}. It is responsible for calculating routes, installing forwarding rules, and enforcing security policies across all underlying domains, enabling complex services like end-to-end network slicing \cite{park2023technology}.
\end{itemize}

By managing the entire network as a single, programmable entity, SDN provides the flexibility and automation necessary for modern services like 5G and the Internet of Things. However, this centralization also underscores the critical need for security. The control plane and the core network components become high-value targets, necessitating the Zero Trust principles and secure framework proposed in this work.



\subsection{Network Operating Systems}
\label{sec:nos}

The NOS is the specialized software that manages the hardware resources of network devices and provides the APIs for implementing control plane functions. For decades, the NOS landscape was dominated by monolithic, proprietary systems like Cisco IOS and Juniper Junos, which tightly coupled the software to specific hardware. The modern shift towards network disaggregation, however, was not a single event but an architectural evolution driven by seminal research in software modularity, performance on commodity hardware, and distributed control.

\paragraph{Modularity in Software Routers}
The first major shift was conceptual, focusing on the software architecture itself. The \textbf{Click Modular Router} pioneered the idea of constructing routers from a collection of fine-grained, reusable packet-processing modules called "elements" \cite{kohler2000click}. In the Click architecture, a router's functionality is defined by a directed graph where packets flow between elements that perform simple tasks like classification, queueing, and header modification. This modularity demonstrated that complex and high-performance packet processing could be achieved with a flexible, extensible software design, laying the intellectual groundwork for the modular, containerized architectures seen in modern NOSs.

\paragraph{Performance on Commodity Hardware}
The second shift addressed the critical question of performance. While modular software was flexible, it was widely believed that only specialized hardware (ASICs) could deliver carrier-grade performance. The \textbf{RouteBricks} project challenged this assumption by demonstrating how to build a high-throughput software router by exploiting parallelism in commodity, multi-core servers \cite{dobrescu2009routebricks}. By parallelizing router functionality both across multiple servers in a cluster and within each server using multi-queue NICs, RouteBricks proved that general-purpose hardware was a viable platform for high-performance networking. This work validated the "whitebox" hardware model that is central to network disaggregation today.

\paragraph{The Distributed Control Platform}
With programmable data planes and high-performance commodity hardware becoming feasible, the final piece was a scalable control plane. The \textbf{Onix} platform was a landmark project that defined the architecture for a distributed "operating system for networks" \cite{koponen2010onix}. Onix introduced the concept of a logically centralized but physically distributed control platform that provides applications with a global view of the network through a \textbf{Network Information Base}. It provided a general API for control logic, abstracting away the complexities of state distribution, failure recovery, and scalability. Onix thus provided the blueprint for modern SDN controller platforms, which are designed to manage a fleet of switches running a device-level NOS.

\paragraph{The Modern Open-Source Landscape}
These foundational concepts - modularity, performance on whitebox hardware, and distributed control - have directly inspired the current generation of open-source SONiC. As the primary focus of this thesis, SONiC is a direct descendant of this evolution. It is highly modular (using a containerized architecture reminiscent of Click's elements), designed for high-performance whitebox switches (as validated by RouteBricks), and intended to be managed by a centralized or distributed controller (as architected by Onix) \cite{alsabeh2020leveraging}.
% \end{itemize}

While these modern NOSs deliver on the promise of a flexible, disaggregated, and high-performance network infrastructure, they primarily focus on functional correctness and performance. They do not, by default, integrate the proactive, identity-driven security mechanisms required for a Zero Trust environment. This thesis addresses that critical gap by designing a framework to extend a modern NOS like SONiC with the security services necessary for deployment in high-threat, carrier-grade networks.

% \subsection{Key Technologies and Ecosystems}
% \subsubsection{The ONAP Orchestration Framework}
% \label{sec:onap_orchestration_framework}

% The framework for orchestration and automation utilized as the foundation for this thesis is the ONAP. ONAP represents a comprehensive, open-source platform designed to automate the lifecycle of network services and functions. Its high-level architecture, depicted in Figure~\ref{fig:onap_architecture}, is fundamentally model-driven and policy-driven, enabling operators to design, instantiate, monitor, and manage complex services with minimal manual intervention \cite{onap_architecture_2024}.

% \begin{figure}[H]
%     \centering
%     \includegraphics[width=\textwidth]{onap_architecture.png}
%     \caption{The ONAP High-Level Architecture illustrating the separation of Design-Time and Run-Time components. Source: ONAP Project.}
%     \label{fig:onap_architecture}
% \end{figure}

% The ONAP architecture is logically divided into two primary environments, which are clearly demarcated in the diagram: a Design-Time framework and a Run-Time framework.

% \paragraph{Design-Time Environment} As shown on the left side of Figure~\ref{fig:onap_architecture}, the Design-Time environment is centered around the \textbf{Service Design \& Creation (SDC)} component. This is where the information models for network resources—such as the YANG model defining the configuration schema for a SONiC switch—are onboarded and packaged into comprehensive service designs. SDC provides the tools to define the resources, their relationships, and the workflows that govern their lifecycle.

% \paragraph{Run-Time Environment} The Run-Time environment, which constitutes the core operational section of the diagram, is responsible for executing the models and workflows defined at design time. While it contains many components, the critical interaction path for this research flows from the \textbf{Service Orchestration (SO)} component to a domain controller. The SO interprets high-level service requests and decomposes them into granular tasks that are delegated to the appropriate controller for execution.

% For this thesis, the most pertinent component is the \textbf{SDN Controller}. This component embodies the principles of SDN (see ~\ref{sec:sdn}) and serves as a key domain controller shown in the lower tier of the Run-Time block. SDN-C's role is to manage the real-time configuration and state of network elements within the \textit{Managed Environment}. As the diagram illustrates, SDN-C acts as the bridge between the abstract orchestration logic and the concrete network device. Critically, its southbound interface was designed to communicate using established, RPC-based protocols, with \textbf{NETCONF} being its primary mechanism.

% This architectural design, specifically the reliance on NETCONF within the SDNC component, creates the central technological gap addressed by this thesis. While SDN-C is a powerful and mature component for managing legacy network devices, it lacks the native capability to interface with modern, cloud-native data planes like SONiC that have standardized on the gNMI protocol. Therefore, the secure framework proposed in this work is designed to act as a specialized, trusted southbound extension for SDNC, enabling it to enforce policy and manage configurations on a SONiC-based network fabric.





% \subsubsection{A Cloud-Native Network Operating System: SONiC}
% \label{sec:sonic_nos}

% Complementing the orchestration layer is the data plane, which in modern disaggregated networks is increasingly powered by open-source, cloud-native network operating systems (NOS). The leading example, and the focus of this thesis's data plane, is the \textbf{SONiC}. Originally developed by Microsoft for its Azure data centers and now a project under the Linux Foundation, SONiC is a fully disaggregated NOS built on Debian Linux, designed to run on a wide variety of whitebox hardware from multiple vendors \cite{SONiC_Wiki}.

% SONiC's architecture fundamentally differs from traditional monolithic network operating systems. Its core design principles are modularity, hardware abstraction, and a database-centric state machine.

% \paragraph{Modular and Containerized Architecture} SONiC encapsulates distinct network functions, such as the BGP routing stack (FRR), LLDP, or SNMP—into individual \textbf{Docker containers} \cite{SONiC_Wiki}. This modularity provides process isolation, simplifies upgrades, and allows for a highly customizable and resilient system.

% \paragraph{Database-Centric Communication} The cornerstone of SONiC's architecture is its reliance on a centralized \textbf{Redis} in-memory database engine. Instead of direct inter-process communication, containers interact by publishing to and subscribing from various tables within the Redis database. Key databases include:
% \begin{itemize}
%     \item \textbf{CONFIG\_DB:} The persistent source of truth for the intended configuration of the switch.
%     \item \textbf{APPL\_DB:} A database where applications write their operational intent (e.g., a new route from the BGP container).
%     \item \textbf{ASIC\_DB:} A hardware-agnostic representation of the state to be programmed into the switching silicon.
% \end{itemize}


% \paragraph{Hardware Abstraction via SAI}
% To achieve true vendor independence, SONiC is built upon the principle of network disaggregation, which decouples the network operating system from the underlying hardware. The key enabler for this disaggregation is the \textbf{SAI}, an Open Compute Project (OCP) standard. SAI provides a standardized, vendor-neutral C-language API to program the forwarding plane of network silicon.

% As illustrated in Figure~\ref{fig:sai}, SAI acts as a stable abstraction layer. The SONiC NOS and its applications are developed against the standard SAI API, without needing any knowledge of the specific ASIC underneath. In turn, each silicon vendor (e.g., Broadcom, Mellanox, Marvell) provides an SAI library, or adapter, that translates the standard SAI API calls into their own proprietary SDK commands.

% \begin{figure}[H]
%     \centering
%     \includegraphics[width=0.8\textwidth]{sai.jpg}
%     \caption{The Disaggregated Network Stack Enabled by SAI. Source: SONiC Project}
%     \label{fig:sai}
% \end{figure}



% This architectural model is what allows a single SONiC software image to run on a multitude of "white box" switch platforms from different hardware manufacturers. Within the SONiC architecture, this interaction is managed by the \textbf{Syncd} container. The workflow is as follows: \texttt{Syncd} subscribes to the hardware-agnostic state in the \texttt{ASIC\_DB}, makes the corresponding standard SAI API calls, and the vendor-provided SAI library executes the necessary hardware-specific instructions to program the ASIC. This model is fundamental to the flexibility of SONiC and ensures that the management framework built in this thesis can operate across a heterogeneous, multi-vendor data plane.

% \begin{figure}[H]
%     \centering
%     \includegraphics[width=\textwidth]{sonic_management_plane.png}
%     \caption{Architecture of a modern NOS management plane. Source: SONiC Project}
%     \label{fig:sonic_management_plane}
% \end{figure}

% This entire architecture is managed through modern, programmatic interfaces. As illustrated in Figure~\ref{fig:sonic_management_plane}, a SONiC device exposes its configuration and operational state through model-driven APIs. There are two protocols that SONiC supports: \textbf{RESTCONF } and \textbf{gNMI }but the native and most efficient protocol for this interaction is the \textbf{gNMI}. A configuration change initiated by a gNMI client follows a precise workflow:
% \begin{enumerate}
%     \item The gNMI server receives a `Set` request and validates it against the device's YANG models.
%     \item Upon successful validation, the server writes the new configuration to the \texttt{CONFIG\_DB}.
%     \item The \textbf{Switch State Service} container, specifically the \textbf{Orchagent} process, detects the change in the database.
%     \item Orchagent translates this high-level configuration into a series of hardware-agnostic instructions and publishes them to the \texttt{ASIC\_DB}.
%     \item The \textbf{Syncd} process consumes these instructions and uses the SAI API to program the physical hardware.
% \end{enumerate}

% This modern, gNMI-native management paradigm is highly efficient and aligns with cloud-native automation principles. However, it is fundamentally incompatible with the NETCONF-based control plane of ONAP's SDN-C. This protocol and paradigm mismatch is precisely what necessitates the development of the secure adapter framework proposed in this thesis, which will serve as the trusted bridge between these two ecosystems.









\section{Security Foundations for Zero Trust}
% \label{sec:}

The dynamic and disaggregated nature of modern networks requires a fundamental shift in security paradigms, moving away from perimeter-based trust to a more rigorous, identity-centric model. This section details the evolution and core principles of Zero Trust Architecture, which forms the security foundation for the adapter developed in this thesis.

\subsection{The Evolution from Perimeter Security to Zero Trust}
Historically, enterprise security was dominated by the \textit{perimeter-based} model, often described as a \textit{castle-and-moat} approach. This model operated on the assumption that all entities inside the internal network-users, devices, and services, were inherently trustworthy, while anything outside was untrustworthy. Consequently, security resources were concentrated on fortifying the network perimeter with firewalls, VPNs, and internet gateways to separate the trusted internal zone from the untrusted external world.



However, this model's foundations began to erode with the rise of modern technology trends such as cloud computing, mobile devices, and widespread remote work, which made the network perimeter increasingly ambiguous. The assumption of implicit internal trust became a significant vulnerability, as insider threats and compromised internal devices grew more common. \cite{mo_perimeter_problem} Major security incidents demonstrated that once an attacker breached the perimeter, they could often move laterally within the network with little resistance, accessing and exfiltrating sensitive data.
 
In response to these challenges, the concept of Zero Trust was formally introduced in 2010 by John Kindervag, an analyst at Forrester Research.\cite{kindervag_2010} He proposed a new philosophy: security should not be based on network location but on a principle of \textit{never trust, always verify}. This paradigm shift dictates that no entity, whether inside or outside the network, should be trusted by default. A landmark implementation of this philosophy was Google's BeyondCorp initiative, which eliminated the traditional corporate VPN and moved to a model where access to services is granted based on verified user and device credentials, regardless of network location.

The model gained widespread recognition and standardization with the publication of the National Institute of Standards and Technology (NIST) Special Publication 800-207, "Zero Trust Architecture," in 2020 \cite{rose_zero_2020}. This sections provides a formal definition and an abstract architecture for implementing Zero Trust, solidifying its role as the leading security strategy for modern enterprises.




\subsection{Zero Trust Architecture}
\label{sec:zta_overview}

Zero Trust is not a single product or technology but a set of cybersecurity paradigms that fundamentally shift defenses from static, network-based perimeters to a focus on users, assets, and resources. As formally defined by the National Institute of Standards and Technology (NIST):

\begin{quote}
    ``Zero trust (ZT) is the term for an evolving set of cybersecurity paradigms that move defenses from static, network- based perimeters to focus on users, assets, and resources. A ZTA uses zero trust principles to plan industrial and enterprise infrastructure and workflows. '' \cite{rose_zero_2020}
\end{quote}

The core philosophy of ZTA is the elimination of implicit trust. In traditional "castle-and-moat" security models, entities inside the network perimeter were often trusted by default. ZTA assumes that the network is always hostile and that external and internal threats exist on the network at all times. Consequently, trust is never granted implicitly but must be continuously evaluated.

\subsubsection{Core Principles of Zero Trust}
\label{sec:zta_principles}

To implement a robust ZTA, organizations must adhere to a set of foundational tenets. Both NIST and the UK National Cyber Security Centre (NCSC) have published guidance to standardize these principles \cite{ncsc_zta} \cite{rose_zero_2020}. The logical progression of a Zero Trust strategy is illustrated in Figure~\ref{fig:zta_principles}, which outlines the seven key tenets required to secure an enterprise environment.

\begin{figure}[H]
    \centering
    \includegraphics[width=0.6\textwidth]{zta_principles.png} % Ensure you upload the image and name it correctly
    \caption{The Seven Tenets of Zero Trust}
    \label{fig:zta_principles}
\end{figure}


\begin{enumerate}
    \item \textbf{Everything is a Resource:} 
    The foundation of ZTA is a comprehensive knowledge of the architecture. This involves identifying all data, computing services, and devices as resources. As noted by the NCSC, knowing the architecture allows organizations to identify key risks and avoid pitfalls when integrating services. Every component, whether a user endpoint or a cloud-native function, must be treated as a resource to be protected.
    
    \item \textbf{Secure All Communications:} 
    Trust must never be placed in the network, including the local network. All communication traffic must be secured regardless of network location. This principle mandates the use of secure transport protocols (such as TLS) to protect data in transit, ensuring that traffic is less susceptible to interception or manipulation by threats residing on the underlying infrastructure.
    
    \item \textbf{Session-Based Access:} 
    Access to individual enterprise resources is granted on a per-session basis. Trust is not transitive; authenticating to the network does not grant access to applications. Each request must be evaluated independently, ensuring that access is granular and time-bound rather than persistent.
    
    \item \textbf{Policies Must Be Dynamic:} 
    Access to resources is determined by dynamic policy, including the observable state of client identity, application, and the requesting asset. The NCSC emphasizes that policies are the power behind ZTA, facilitating risk-managed sharing of data. These policies must adapt to the changing risk context of the environment rather than relying on static rules.
    
    \item \textbf{Monitor Security Posture:} 
    The enterprise must monitor and measure the integrity and security posture of all owned and associated assets. No asset is inherently trusted. The NCSC guidance highlights that monitoring should focus on user behavior, device health, and service status to establish confidence signals for the policy engine.
    
    \item \textbf{Authenticate Before Connect:} 
    All resource authentication and authorization are dynamic and strictly enforced before access is allowed. This principle, often summarized as \textit{authenticate and authorize everywhere,} requires that the system verify the user identity, device health, and other attributes before a connection to the data plane is ever established.
    
    \item \textbf{Measure and Improve:} 
    The enterprise collects as much information as possible about the current state of assets, network infrastructure, and communications and uses it to improve its security posture. This creates a feedback loop where monitoring data feeds back into policy refinement, ensuring the architecture evolves to meet new threats.
\end{enumerate}

In the context of this thesis, the proposed Secure Adapter is designed to enforce specifically principles 2, 4, and 6: it secures communications via mTLS, enforces dynamic policies via RBAC, and strictly authenticates the orchestration controller before allowing any connection to the SONiC data plane.


\subsubsection{ZTA Logical Architecture}
\label{sec:zta_logical_architecture}

To understand how the Secure Framework enforces Zero Trust, it is essential to map its components to the logical architecture defined by the National Institute of Standards and Technology (NIST) in its seminal publication, SP 800-207 \cite{rose_zero_2020}. This standard provides an abstract, vendor-neutral model for ZTA, which is depicted in Figure~\ref{fig:nist_zta_model}.

\begin{figure}[H]
    \centering
    \includegraphics[width=0.8\textwidth]{nist_zta_model.png} 
    \caption{Core Zero Trust Logical Components. Adapted from NIST SP 800-207 \cite{rose_zero_2020}.}
    \label{fig:nist_zta_model}
\end{figure}

The core of the NIST model is the interaction between a \textbf{Policy Engine (PE)}, a \textbf{Policy Administrator (PA)}, and a \textbf{PEP}.
\begin{itemize}
    \item The \textbf{Policy Engine} is the decision-making brain of the architecture. It evaluates access requests by analyzing multiple inputs including enterprise policies, threat intelligence data, device compliance status, activity logs, and other contextual signals. Based on this comprehensive analysis, the PE determines whether to grant or deny access to requested resources.
    \item The \textbf{PA} acts as the communication bridge between the Control Plane and Data Plane. The PA receives access decisions from the PE and communicates these decisions to the Policy Enforcement Point. It generates session-specific authentication tokens or credentials that enable the PEP to enforce the PE's decisions.
    \item The \textbf{PEP} is the active enforcement component that sits in the communication path between subjects (users/devices) and enterprise resources.
\end{itemize}


By instantiating these logical ZTA components, the Secure Framework moves beyond simple security and becomes an active participant in the dynamic, policy-driven access control required for a true Zero Trust environment.

\subsubsection{Trusted Algorithm}
\label{sec:trusted_algorithm}

Within a Zero Trust Architecture, the Policy Engine (PE) serves as the central \textit{brain} of the system, and the \textbf{Trust Algorithm (TA)} represents its primary thought process. As defined in NIST SP 800-207, the Trust Algorithm is the logic used to ultimately grant or deny access to a resource. It moves beyond static, rule-based decisions by dynamically aggregating and weighting inputs from multiple data sources to calculate a confidence level for each access request \cite{rose_zero_2020}.

The inputs required for the Trust Algorithm to make an informed decision are categorized into five distinct groups, as illustrated in Figure~\ref{fig:trust_algorithm}.

\begin{figure}[t]
    \centering
    \includegraphics[width=0.8\textwidth]{trust_algorithm.png} % Ensure you upload the image and name it trust_algorithm.png
    \caption{Trust Algorithm Input. Source: Adapted from NIST SP 800-207 \cite{rose_zero_2020}.}
    \label{fig:trust_algorithm}
\end{figure}


\begin{itemize}
    \item \textbf{Access Request:} This encompasses the immediate context of the request, including the identity of the requester and the specific resource being targeted. It also includes ephemeral attributes such as the operating system version, the specific application making the request, and its patch level.
    
    \item \textbf{Subject Database:} This source provides the "who" of the transaction. It references an identity management system (such as LDAP or PKI) to retrieve the subject's attributes, assigned privileges, and historical behavior patterns.
    
    \item \textbf{Asset Database:} This database maintains the known state of enterprise assets. The TA compares the observable status of the requesting asset (e.g., location, integrity, software presence) against this known baseline to detect anomalies.
    
    \item \textbf{Resource Policy Requirements:} This defines the minimal security requirements for the specific resource being accessed. Policies may dictate required authenticator assurance levels (e.g., Multi-Factor Authentication), geographic restrictions, or data sensitivity classifications.
    
    \item \textbf{Threat Intelligence:} This involves external or internal feeds regarding active threats, malware signatures, and command-and-control node identifiers. This dynamic input allows the TA to deny access if the requesting asset is exhibiting behavior associated with known attacks, even if it has valid credentials.
\end{itemize}

The Trust Algorithm processes these inputs, applying proprietary or configurable weights to each factor to derive a final decision. This decision is then passed to the \textbf{PA} for execution. The PA is responsible for configuring the PEP to either establish the secure communication path or terminate the connection based on the TA's verdict.


\subsection{Zero Trust Network Access}
\label{sec:ztna}

While ZTA provides the overarching strategy, ZTNA represents its practical application to securing access to private applications and resources. 
This model inverts the traditional security paradigm by making applications invisible to the public internet and treating the internal network as untrusted. Access is brokered on a per-session basis, strictly enforcing the principles of least privilege and prohibiting lateral movement.

\subsubsection{History and Evolution of Private Access}
\label{sec:ztna_history}

The concept of ZTNA did not emerge in a vacuum; it is the culmination of a decades-long evolution in secure networking, with each preceding technology addressing the limitations of its predecessor.

\paragraph{The Perimeter Era: IPSec and VPNs (1995-1998)}
The dominant paradigm for remote access for over two decades was the VPN, often built upon the IP Security (IPSec) protocol suite \cite{rfc2401}. The goal of a VPN was to create a secure, encrypted tunnel over the public internet, effectively extending the trusted corporate perimeter to a remote user or a branch office. Once authenticated, the user's device was placed "on" the internal network, granting broad access to network resources. This model is the antithesis of Zero Trust, as it operates on a binary "trust vs. untrust" model and, once breached, provides an ideal environment for an attacker's lateral movement.

\paragraph{The Internal Control Era: NAC (2003)}
A significant step towards an identity-centric model was the introduction of Network Access Control (NAC), standardized as IEEE 802.1X \cite{ieee8021x}. NAC introduced a crucial concept: authenticating a user or device \textit{before} granting any access to the local area network (LAN). By acting as a gatekeeper at the network edge, NAC provided protection against unauthorized devices connecting to an otherwise open internal network. While still perimeter-focused, NAC was a critical precursor to ZTNA as it was one of the first widely adopted technologies to tie network access directly to a verified identity.

\paragraph{The Precursors to ZTNA: Software-Defined Perimeter (2007)}
The direct architectural ancestor of ZTNA is the Software-Defined Perimeter (SDP), a framework developed from work at the U.S. Defense Information Systems Agency (DISA). The SDP model, also known as "Black Cloud," implements the core tenets of ZTNA:
\begin{enumerate}
    \item The application infrastructure is made \textit{dark} to the network, with no listening ports visible to unauthorized users.
    \item A "Single Packet Authorization" (SPA) is sent to a controller, which authenticates the user and device.
    \item If authorized, the controller dynamically provisions a secure, encrypted, per-session connection directly between the user's device and the specific application.
\end{enumerate}
This model, which heavily influenced Google's BeyondCorp initiative (2014) \cite{BeyondCorp2014}, effectively creates a dynamic, individualized perimeter around each application for each user, for each session.

\paragraph{Modern ZTNA Implementations: Mesh Overlays (2018-Present)}
The latest evolution of ZTNA leverages agent-based mesh overlay networks. In this model, a central trust broker manages identity and policy, but it does not sit in the data path. Instead, authenticated devices establish direct, end-to-end encrypted tunnels to each other using modern cryptographic protocols like WireGuard \cite{donenfeld2017wireguard}. These connections are typically established using outbound-only traffic and NAT traversal techniques, allowing firewalls to remain completely closed to inbound connections. This dramatically reduces the public attack surface and provides the most resilient and performant implementation of ZTNA principles to date, enabling both north-south (user-to-application) and east-west (application-to-application) secure communication.




% \subsubsection{Software-Defined Perimeter}
% \label{sec:sdp}

% The Software-Defined Perimeter (SDP) is a security framework that implements the principles of Zero Trust to protect application infrastructure. It inverts the traditional perimeter security model by making resources invisible (or "black") to all unauthorized entities. The Cloud Security Alliance (CSA), which standardized the architecture, defines SDP as a framework that provides an on-demand, dynamically provisioned, and trusted overlay network to mitigate network-based attacks \cite{csa_sdp_spec_v2}.

% The core concept of SDP is a strict "need-to-know" model, where a device's identity must be verified and authenticated \textit{before} it is granted any access to the application infrastructure \cite{moubayed2019sdp}. This approach effectively prevents a wide range of network attacks, including port scanning, credential theft, and Denial of Service (DoS), by eliminating the public attack surface.

% \paragraph{SDP Architecture}
% The SDP architecture, which evolved from the U.S. Department of Defense's "Black Core" network initiative, is logically separated into a control plane and a data plane. It consists of three primary components, as shown in Figure~\ref{fig:sdp_architecture}:

% \begin{figure}[H]
%     \centering
%     \includegraphics[width=0.9\textwidth]{sdp_framework.png} % Ensure you upload the image and name it correctly
%     \caption{The Software-Defined Perimeter Framework Architecture. Source: Adapted from \cite{moubayed2019sdp}.}
%     \label{fig:sdp_architecture}
% \end{figure}

% \begin{itemize}
%     \item \textbf{SDP Controller:} This is the "brain" of the architecture, acting as the trust broker or Policy Decision Point (PDP). It integrates with identity providers (e.g., IAM, MFA) to authenticate users and devices, evaluates access policies, and communicates authorizations to the other components.
%     \item \textbf{Initiating Host:} This is the client device (e.g., a user's laptop or a server) that is requesting access to a protected resource.
%     \item \textbf{Accepting Host:} This is the server or gateway that hosts the protected application or service. By default, the AH is configured to drop all traffic from any source except the trusted SDP Controller.
% \end{itemize}

% \paragraph{SDP Workflow}
% The SDP workflow strictly enforces an "authenticate before connect" sequence, which is its key security advantage. The process involves two distinct phases:

% \begin{enumerate}
%     \item \textbf{Control Plane Authentication via Single Packet Authorization (SPA):} The Initiating Host (IH) does not attempt to connect directly to the application. Instead, it sends a specially crafted, encrypted, and authenticated message known as a Single Packet Authorization (SPA) to the SDP Controller. The SPA packet contains the user's credentials and the requested resource. Because the Accepting Host (AH) has no open ports, it is immune to SPA-based DoS attacks and port scans. The Controller validates the SPA, and if the user and device are authorized, it instructs the AH to temporarily accept a connection from that specific IH.

%     \item \textbf{Data Plane Connection via mTLS):} Once authorized by the Controller, the IH is given the necessary information to connect directly to the AH. This connection is established as a dynamic, per-session, end-to-end encrypted tunnel using mTLS. This ensures that both the client and the server are cryptographically authenticated to each other, and all subsequent communication is confidential and integral.
% \end{enumerate}

% By separating the authentication/authorization (control plane) from the application access (data plane), SDP creates a highly secure overlay network that is resilient to network-level attacks. The principles of SDP—specifically the "authenticate before connect" model and the use of a trusted broker to dynamically provision access—are directly applicable to the Secure Framework proposed in this thesis, which applies these same concepts to secure the network management plane.

% \subsubsection{Overlay Networks}
% \label{sec:overlay_networks}

% Overlay networks are a powerful networking paradigm that involves layering a logical, virtual network on top of an existing physical network infrastructure, known as the underlay. A key advantage of this approach is that it allows for the incremental deployment of new services and functionalities—such as virtual networking, enhanced QoS, or improved routing—without requiring a disruptive and costly upgrade of the entire underlying physical infrastructure \cite{paillisse2021next}.

% The reference architecture for an overlay network, as detailed in \cite{paillisse2021next}, consists of four distinct components that work in concert, as illustrated in Figure~\ref{fig:overlay_architecture}.

% \begin{figure}[h!]
%     \centering
%     \includegraphics[width=\textwidth]{overlay_architecture.png} % Ensure you upload the image and name it correctly
%     \caption{Reference Architecture for Overlay Networks. Source: Adapted from \cite{paillisse2021next}.}
%     \label{fig:overlay_architecture}
% \end{figure}

% \begin{itemize}
%     \item \textbf{Underlay Network:} This is the physical, underlying network layer, typically a standard IP network, composed of routers and switches that provide connectivity between the overlay's edge devices.
    
%     \item \textbf{Overlay Network:} This is the logical, upper network layer. Its primary function is performed by \textbf{en/de-capsulation routers} located at the network edge. These routers are responsible for adding (encapsulating) and removing (decapsulating) the overlay headers from packets.
    
%     \item \textbf{Overlay Endpoints:} These are the end hosts (servers, computers, devices) that communicate with each other. A crucial characteristic of the overlay model is that the endpoints are unaware of the underlay's existence; they communicate as if they were on a simple, direct network.
    
%     \item \textbf{Mapping System:} This is a logically centralized database that maintains the mappings between overlay identifiers (e.g., the logical IP address of an endpoint) and underlay locators (e.g., the physical IP address of the en/de-capsulation router responsible for that endpoint). As noted in \cite{paillisse2021next}, this component is sometimes optional (e.g., in flood-and-learn VXLAN deployments) and can range from a simple key-value store to a fully-fledged SDN controller.
% \end{itemize}

% The fundamental mechanism of an overlay network is \textbf{encapsulation}. As shown in the packet format diagram in Figure~\ref{fig:overlay_architecture}, an original packet from an endpoint (Payload) is first wrapped with an \textit{Overlay IP header}. This entire packet is then treated as the payload for the underlay network and is wrapped with an \textit{Underlay IP header}. This outer header is used to route the packet across the physical infrastructure from the source encapsulation router to the destination decapsulation router, where the headers are stripped in reverse order and the original packet is delivered to the final endpoint.

% The concept of overlay networks is particularly relevant to modern security architectures. Frameworks like SDP and other ZTNA solutions are, in fact, specialized security overlays. They leverage this model to create secure, encrypted, and logically isolated perimeters on top of untrusted underlay networks like the public internet.

\subsubsection{OpenZiti}
\label{sec:openziti}

OpenZiti is a prominent open-source project that exemplifies the principle of application-embedded Zero Trust networking. Rather than focusing solely on securing the network layer, OpenZiti provides a complete framework, including Software Development Kits (SDKs), to integrate Zero Trust principles directly into an application's communication logic \cite{openziti_project}. This approach shifts the security perimeter from the network edge to the application identity itself, enabling secure connectivity over any underlying network, including the public internet, without reliance on traditional VPNs.

\paragraph{Architectural Overview}
The OpenZiti architecture, illustrated in Figure~\ref{fig:openziti_architecture}, is a full-featured overlay network that strictly separates its control and data planes.

\begin{figure}[t]
    \centering
    \includegraphics[width=\textwidth]{openziti_architecture.png} % Ensure you upload the image and name it correctly
    \caption{The OpenZiti Zero Trust Overlay Network Architecture. Source: OpenZiti Project}
    \label{fig:openziti_architecture}
\end{figure}

\begin{itemize}
    \item \textbf{Controller:} This is the central brain of the network, acting as the trust broker and PDP. As shown by the control message flows in the diagram, it is responsible for authenticating and authorizing all network participants, managing policies, and orchestrating the establishment of secure communication paths.
    
    \item \textbf{Overlay Network Fabric:} The data plane is composed of a mesh of \textbf{Edge Routers} and internal fabric routers. These routers do not trust any traffic by default. They only forward packets based on the secure sessions authorized and set up by the Controller.
    
    \item \textbf{Endpoints (SDKs and Tunnelers):} An application can join the OpenZiti network in one of two ways. The preferred method is by embedding an \textbf{OpenZiti SDK}, which makes the application a first-class citizen of the overlay. For legacy applications that cannot be modified, OpenZiti provides tunneling applications that act as local proxies, intercepting traffic and forwarding it securely over the overlay.
\end{itemize}

\paragraph{Zero Trust Workflow}
The OpenZiti workflow strictly adheres to the \textit{authorize before connect} principle. An application with an embedded SDK first establishes a secure control channel to the Controller to authenticate itself. The Controller verifies the identity and evaluates its policies to determine which services the application is authorized to access. If authorized, the Controller instructs the relevant Edge Routers to permit a connection from the client to the target service. The client SDK then establishes a secure, end-to-end encrypted data path through the overlay fabric to the service.

As depicted on the right side of Figure~\ref{fig:openziti_architecture}, direct access to the target API or service is blocked. The service is effectively \textit{dark} to the network, and the only way to establish a connection is by being explicitly authorized by the Controller. This model provides least-privilege access and prevents lateral movement, as a compromised client can only access the specific services it was granted, and only for the duration of the authorized session.


\subsection{Enabling Technologies for the Secure Framework}
\label{sec:enabling_technologies}

The realization of the ZTA logical components within the Secure Framework is made possible by the integration of two key, industry-standard technologies: mTLS for authentication andRBAC for authorization.

\subsubsection{Mutual TLS for Strong Authentication}
\label{sec:mtls}
Standard TLS provides one-way authentication: the client verifies the identity of the server to prevent man-in-the-middle attacks. However, in a Zero Trust environment, this is insufficient. The server (the resource) must also be able to cryptographically verify the identity of the client (the subject). This is the role of mTLS.

As standardized in RFC 8446 for TLS 1.3, the mTLS handshake extends the standard TLS process by requiring the client to also present an X.509 certificate to the server \cite{rescorla_tls_2018}. Both parties then validate each other's certificates against a trusted Certificate Authority. This process establishes a secure, encrypted channel where the identities of both endpoints are unequivocally proven.

Within our framework, mTLS is the cornerstone of the \textit{never trust, always verify} principle. It is used on both the northbound (SDN-C to Framework) and southbound (Framework to SONiC) connections. This ensures that every component in the management chain is strongly authenticated, eliminating the possibility of unauthorized or spoofed entities issuing configuration commands. The authenticated identity extracted from the client's certificate becomes the trusted principal for the subsequent authorization step.

\subsubsection{Role-Based Access Control for Granular Authorization}
\label{sec:rbac}
Once a subject's identity has been authenticated via mTLS, the next step in the Zero Trust workflow is to determine what that subject is \textit{authorized} to do. For this, the framework implements a RBAC model. RBAC is a widely adopted policy model that simplifies permission management by assigning permissions to roles rather than directly to individual users or services \cite{ferraiolo_proposed_1992}.

The RBAC model within the Secure Framework consists of three core elements:
\begin{enumerate}
    \item \textbf{Principals:} The authenticated identities of the clients, derived from their mTLS certificates (e.g., `sdnc.onap`).
    \item \textbf{Roles:} A named set of permissions (e.g., `interface-config-admin`, `read-only-operator`). A principal is mapped to one or more roles.
    \item \textbf{Permissions:} A specific rule that grants access to perform an operation on a resource. In our framework, a permission is a tuple defining the allowed `(operation, admin)`.
\end{enumerate}

When a request arrives, the framework's internal Policy Engine retrieves the roles associated with the authenticated principal. It then checks if any of these roles contain a permission that explicitly allows the requested operation on the target resource. If and only if such a permission exists, the request is authorized. This mechanism provides the fine-grained, policy-driven control necessary to enforce the principle of least privilege, ensuring that even a trusted controller can only perform the exact actions it has been explicitly permitted to.

\section{Network Management Protocols}
\label{sec:net_mgmt_protocols}

The evolution from monolithic to disaggregated networks has been paralleled by an evolution in the protocols used to manage them. The choice of management protocol fundamentally dictates how an orchestrator interacts with network elements, influencing performance, scalability, and operational paradigms. This section details the two distinct models at the heart of this thesis: the traditional RPC-based approach embodied by NETCONF and the modern, streaming-first approach of gNMI.

\subsection{The RPC-Based Model: NETCONF and YANG}
\label{sec:rpc_model}

The Remote Procedure Call (RPC)-based model represents a mature, robust paradigm for network configuration, heavily influenced by the requirements of carrier-grade networking. This model is formalized through the combination of the NETCONF protocol and the YANG data modeling language.

\paragraph{NETCONF}
The Network Configuration Protocol (NETCONF) provides a standardized mechanism for installing, manipulating, and deleting the configuration of network devices \cite{rfc6241_netconf}. It was designed to overcome the limitations of legacy methods like Command-Line Interface (CLI) scraping and Simple Network Management Protocol (SNMP), which were often vendor-specific, error-prone, and lacked transactional integrity. Key characteristics of NETCONF include:
\begin{itemize}
    \item \textbf{Transactional Integrity:} NETCONF operations are transactional, allowing a set of configuration changes to be applied in an all-or-nothing manner. This prevents devices from being left in an inconsistent state.
    \item \textbf{Separation of Datastores:} It defines distinct configuration datastores, such as \texttt{<running>}, \texttt{<candidate>}, and \texttt{<startup>}, which enables robust configuration workflows (e.g., editing a candidate configuration before committing it to the running state).
    \item \textbf{Structured Data:} All operations and data are encoded in XML over a secure and connection-oriented transport, typically SSH.
    \item \textbf{Standardized Operations:} It provides a small set of base protocol operations, including \texttt{<get-config>}, \texttt{<edit-config>}, and \texttt{<copy-config>}.
\end{itemize}

\paragraph{YANG}
While NETCONF defines the operational framework, YANG (Yet Another Next Generation) provides the data modeling language used to define the structure of configuration and state data \cite{rfc7950_yang}. A YANG model is a formal, human-readable contract that describes the complete operational schema of a device, including data hierarchies, types, constraints, and semantics. This model-driven approach, where the data schema is decoupled from the protocol, is a cornerstone of modern network automation. These include the ability to program the data plane in a protocol-independent manner for adaptive monitoring and control policy enforcement, real-time processing of streaming monitoring data, predictive machine learning for closed-loop network management, and orchestration algorithms for cost-effective, resilient, and efficient service provisioning \cite{clemm2020network}. It also enables automated validation of configurations and provides a stable, programmatic interface for controllers like ONAP's SDN-C.

Together, NETCONF and YANG provide a powerful, transactional framework for precise network configuration. This robustness has made it the standard for traditional telecommunication systems and is the native southbound protocol for key ONAP components.

\subsection{Streaming Model: gRPC Network Management Interface}
\label{sec:streaming_model}

In contrast to the RPC-based model, the streaming model was born from the operational needs of hyperscale data centers, where high-frequency telemetry and efficient, scalable configuration are paramount \cite{singh2015jupiter}. The gNMI is the leading protocol embodying this paradigm.

\paragraph{Foundations: gRPC and Protocol Buffers}
gNMI is built on gRPC (gRPC Remote Procedure Calls), a high-performance RPC framework from Google. Unlike NETCONF's use of XML over SSH, gRPC leverages HTTP/2 for transport, enabling features like bidirectional streaming and header compression. For data serialization, it uses Protocol Buffers (Protobuf), a binary format that is significantly more compact and efficient to parse than XML.

\paragraph{gNMI Operations}
gNMI unifies configuration management and telemetry into a single, elegant protocol with three primary RPCs:
\begin{itemize}
    \item \textbf{Get:} Retrieves a snapshot of data from the target device, similar to NETCONF's \texttt{<get-config>}.
    \item \textbf{Set:} Modifies the state of the target device, supporting transactional updates, replacements, and deletions, analogous to \texttt{<edit-config>}.
    \item \textbf{Subscribe:} This is gNMI's defining feature. It allows a client to request a long-lived stream of data from the target. The client can specify which data it is interested in (via paths) and how it wants to receive updates (e.g., periodically or on-change). This capability is essential for modern, real-time network monitoring and closed-loop automation, representing a significant advancement over the poll-based nature of SNMP or NETCONF notifications \cite{boutaba2018survey}.
\end{itemize}

While gNMI can use YANG as its modeling language, the on-the-wire protocol is fundamentally different. Its efficiency, scalability, and native support for streaming telemetry have made it the de facto standard for cloud-native network operating systems like SONiC. This adoption creates the critical protocol mismatch with legacy orchestration systems that this thesis directly addresses.



\section{Chapter Summary}
This chapter has laid the groundwork by introducing the key technologies and security principles relevant to this thesis. We have established the context of network disaggregation and introduced the primary actors: the ONAP orchestrator and the SONiC data plane. We have detailed the protocol mismatch between NETCONF and gNMI that necessitates a translation layer. Finally, we have defined the core security concepts - ZTA, ZTNA, mTLS, and RBAC - that must be integrated into this layer to build a secure and trustworthy automation framework. The following chapter will review existing literature to identify the research gap that this work addresses.
% Fill the chapter summary


\newpage
\chapter{Related Work}
\label{chap:related}

This chapter surveys the existing body of research in three key areas that form the foundation of this thesis: network orchestration and automation, cross-protocol communication in software-defined networks, and security for disaggregated environments. By examining the state-of-the-art in each domain, we identify the critical technological and security gaps that motivate the design of a secure, policy-driven framework for Zero Trust enforcement.

\section{Network Orchestration and Automation}
\label{sec:related_orchestration}
The paradigm of network orchestration has matured significantly, driven by the architectural principles laid out by the European Telecommunications Standards Institute's (ETSI) Network Functions Virtualization (NFV) framework \cite{han2015network}. The ETSI Management and Orchestration (MANO) architecture defines a three-tiered structure comprising a high-level NFV Orchestrator (NFVO), a Virtualized Network Function Manager (VNFM), and a Virtualized Infrastructure Manager (VIM). This centralized management system is a critical component not only for coordinating programmable networks but also for ensuring network security

The ONAP stands as one of the most comprehensive implementations of the ETSI MANO vision. Research on ONAP has extensively covered its use in automating the lifecycle of complex network services, such as 5G network slicing and SD-WAN deployments. These works highlight ONAP's strength in high-level, intent-based service design and its reliance on domain-specific controllers, like SDN-C, to interact with the underlying network infrastructure. Similarly, Open Source MANO (OSM) is another prominent ETSI-aligned orchestrator, often noted for its lightweight nature and interoperability with VIMs like OpenStack and Kubernetes \cite{Mamushiane2019}.

Beneath the orchestration layer, various open-source projects have emerged to handle the control plane and specific network domains. Notable examples include:
\begin{itemize}
\item \textbf{ONOS (Open Network Operating System):} Developed by ON.Lab, this is a carrier-grade SDN operating system designed for high availability and scalability in service provider and core networks.
\item \textbf{OpenDaylight (ODL):} Supported by the Linux Foundation, ODL provides a modular SDN framework that facilitates the development of network applications and remote device management.
\item \textbf{Faucet:} A lightweight open-source SDN controller based on OpenFlow, designed for enterprise and campus networks with a focus on ease of deployment.
\item \textbf{Tungsten Fabric:} Formerly OpenContrail, this project specializes in SDN and NFV, supporting the deployment of virtual network services.
\item \textbf{CORD (Central Office Re-architected as a Datacenter):} An ON.Lab initiative focused on transforming telecommunications central offices into cloud-native data centers using SDN and NFV technologies.
\end{itemize}

While these platforms provide powerful capabilities, the literature consistently shows that the interaction between high-level orchestrators (like ONAP or OSM) and the data plane is often abstracted. They delegate device-specific configuration tasks to southbound controllers (like ODL or ONOS), assuming the existence of a functional and secure communication channel. However, the internal mechanisms of these interactions in heterogeneous environments are often treated as a solved problem.

\section{Cross-Protocol Communication in SDN}
\label{sec:related_cross_protocol}
The challenge of integrating network devices that speak different management protocols is a long-standing issue in SDN. In early SDN deployments, significant research was dedicated to creating "hybrid" environments where OpenFlow-based controllers could manage legacy devices supporting only SNMP or proprietary CLIs. This often involved creating protocol gateways or abstraction layers that translated controller intent into device-specific commands. While demonstrating feasibility, these initial solutions were often bespoke, difficult to maintain, and lacked a standardized, model-driven approach.


As the field matured, the focus shifted towards model-driven translation, particularly between IETF-standardized protocols like NETCONF/YANG and RESTCONF. The authors in \cite{Alghamdi2022} explore the design of unified northbound interfaces for SDN controllers that can abstract away the specific southbound protocol being used. These layers function by mapping a generic API call to either a NETCONF RPC or a RESTCONF request based on the target device's capabilities, leveraging the common YANG data model as the source of truth.

More recently, with the rise of gNMI in cloud-native data centers, the specific challenge of mediating between NETCONF and gNMI has been addressed by the open-source community. Projects such as the OpenConfig \textbf{gNMI-NETCONF Adapter} and various Python-based libraries (e.g., `pygnmi-netconf`) have emerged. These tools act as \textbf{functional gateways} or protocol mediation layers. Their primary objective is to provide a translation service, allowing a NETCONF-based client (like an ONAP controller) to configure a gNMI-enabled device (like a SONiC switch) by converting NETCONF RPCs into gNMI `Set` requests and vice-versa.


However, a critical analysis of these existing tools reveals a profound limitation in the context of modern security architectures. Their design philosophy is centered on achieving functional protocol correctness, not on serving as a robust security enforcement point. The security model of these adapters is typically confined to the transport layer—for instance, terminating a client's SSH session and initiating a new TLS session to the target device. While this secures the communication channels, it does not provide the granular, identity-aware authorization required for a Zero Trust Architecture. These gateways generally lack an integrated RBAC engine capable of inspecting the content of a management request and authorizing the specific operation against a policy. They are architected to be transparent translators, not opinionated security intermediaries.

Therefore, while the problem of basic protocol translation between NETCONF and gNMI has been addressed by existing open-source tools, they do not solve the critical problem of secure, policy-driven orchestration. They answer the question "Can we translate the request?" but not the more important question, "Even if we can translate it, \textit{should} this specific operation be allowed?" This distinction forms a cornerstone of the research gap that this thesis addresses by designing a framework that is not merely a translator, but a secure Zero Trust Policy Enforcement Point.



% \section{Security in Disaggregated and Virtualized Networks}
% \label{sec:related_security}
% The shift to disaggregated and virtualized networks, particularly those managed by SDN, fundamentally dismantles the traditional perimeter-based security model. The centralization of control and the exposure of the southbound interface (SBI) create new, high-value attack surfaces. Consequently, a significant body of research has focused on developing novel security mechanisms beyond simple transport-layer encryption. These approaches can be broadly categorized into proactive defenses that change the attack surface and advanced cryptographic methods for securing the control channel.

% \section{Proactive Defense: Moving Target Defense}
% Moving Target Defense (MTD) has emerged as a leading proactive defense strategy. Instead of waiting to react to an attack, MTD aims to continuously change the system's attack surface, thereby invalidating information an attacker has gathered during reconnaissance \cite{cho2020survey}. In the context of SDN, a common MTD technique is IP address shuffling, where the SDN controller periodically and randomly changes the virtual IP addresses of hosts. As demonstrated by Mendonça et al., this technique can significantly increase the "Mean Time to Compromise" (MTC) by forcing an attacker to constantly re-scan the network to find their target \cite{mendonca2022security}. While highly effective at thwarting reconnaissance, MTD's primary focus is on obfuscation—making the target difficult to find. It does not, by itself, prevent an attack from an adversary who already possesses valid credentials or has an established session.

% \paragraph{Securing the Control Channel}
% The communication channel between the SDN controller and the data plane devices (the southbound interface) is a critical point of vulnerability. While TLS is the standard recommendation, its reliance on complex Public Key Infrastructure (PKI) management and its susceptibility to various protocol attacks have led researchers to explore alternatives.
% \begin{itemize}
%     \item \textbf{Identity-Based Cryptography (IBC):} Lam et al. propose securing both southbound and east/west-bound SDN communication using IBC \cite{lam2015securing}. IBC simplifies key management by allowing public keys to be derived directly from an entity's identity (e.g., its MAC address or domain name), eliminating the need for a complex certificate lifecycle management system. This approach provides strong channel security with lower operational overhead than traditional TLS.
    
%     \item \textbf{Blockchain Technology:} To address the single point of trust in a traditional Certificate Authority, Ahmad and Mir propose a decentralized, blockchain-based security mechanism for the southbound interface \cite{ahmad2023securing}. In this model, a distributed ledger is used to manage the identities and public keys of controllers and switches. This provides a tamper-proof, auditable record of all authenticated entities in the network, enhancing resilience against attacks targeting a central trust anchor.
% \end{itemize}
% Both IBC and Blockchain offer significant improvements for authenticating the \textit{endpoints} of the control channel. However, their focus remains on securing the session. They ensure that a legitimate controller is talking to a legitimate switch, but they do not provide a mechanism to inspect and authorize the \textit{content} of the commands being sent over that secure channel.


\section{Related Work}
\label{sec:related}


The increasing disaggregation of network functions has spurred research into applying Zero-Trust (ZT) principles to specific architectural domains like SDN and 6G. For instance, Erel-Özçevik~\cite{Erel-Ozcevik2024} proposes a ``Token as a Service'' (TaaS) framework for multi-tenant SDN environments, using JSON Web Tokens (JWT) to manage controller authentication. While this effectively reduces the load on centralized authentication servers, it primarily addresses the identity of the controllers themselves, not the fine-grained authorization of the management commands they issue. Similarly, Duy et al.~\cite{duy2022} propose B-DAC, a blockchain-based framework to secure the SDN northbound interface, ensuring tamper-proof and decentralized access control for applications. These solutions, along with the 6G RAN framework by Sedjelmaci et al.~\cite{Sedjelmaci2024}, exemplify domain-specific ZT solutions that secure vertical technology stacks. Lam et al.~\cite{lam2015securing} made a valuable contribution by using Identity-Based Cryptography (IBC) to secure communication channels in distributed SDN. 


A significant body of research has explored proactive defense paradigms and novel access control mechanisms. One prominent approach is Moving Target Defense (MTD), which aims to confuse attackers by continuously altering the attack surface. As detailed in surveys by Cho et al.~\cite{cho2020survey} and applied to SDN by Mendonça et al.~\cite{mendonca2022security}, MTD enhances resilience through obfuscation. This philosophy differs fundamentally from ZT, which relies on explicit, identity-based verification for every transaction. Other research has focused on securing the data plane itself. Musumeci et al.~\cite{Musumeci2021}, for example, leverage P4-programmable switches for ML-based DDoS attack detection and  Grewal et al.~\cite{grewal2022p4bid} use an information-flow control type system within P4 to enforce security properties at the language level, demonstrating the power of in-network enforcement for data-plane threats. While powerful for data-plane protection, these approaches do not address the distinct challenges of securing the management plane. To establish trust in decentralized systems, Ding and Sato~\cite{Ding2022} proposed ``Bloccess,'' a framework using blockchain for tamper-proof access control. While offering robust, decentralized trust, this approach introduces scalability and latency trade-offs.
% \section{Research Gaps}
% \label{sec:research_gap}
% The review of state-of-the-art security mechanisms for SDN reveals a clear and critical research gap. While existing solutions provide powerful tools, they address different facets of the security problem in isolation. The gap emerges from the limitations of these specialized approaches when considered in the context of a holistic, end-to-end security architecture, as summarized in Table~\ref{tab:approach_summary}.

% \begin{table}[H]
% \centering
% \caption{Summary and Limitations of Existing Security Approaches}
% \label{tab:approach_summary}
% \begin{tabular}{@{}>{\raggedright\arraybackslash}p{3.5cm}p{4.5cm}p{6cm}@{}}
% \toprule
% \textbf{Approach} & \textbf{Primary Focus} & \textbf{Architectural Limitation} \\
% \midrule
% Moving Target Defense (MTD) & 
% Obfuscating the attack surface (e.g., IP shuffling) & 
% Does not authorize actions within an established session. Focuses on \textit{finding} the target, not \textit{accessing} it. \\
% \addlinespace
% Channel Security (TLS, IBC, Blockchain) & 
% Authenticating endpoints (Controller $\leftrightarrow$ Switch) & 
% Does not authorize the \textit{content} of the communication. Secures the session, not the specific command. \\
% \bottomrule
% \end{tabular}
% \end{table}

% The research gap can be defined by two key points:
% \begin{enumerate}
%     \item \textbf{Lack of Command-Level Authorization:} Proactive defenses like MTD and advanced channel security protocols like IBC and Blockchain are designed to either hide the target or authenticate the communicating parties. They do not provide a native mechanism for granular, command-level authorization. They cannot answer the question: "Is this authenticated controller permitted to execute this specific `edit-config` operation on this particular switch?"
    
%     \item \textbf{Absence of a multi-layered security model:} There is no existing framework that combines a secure, identity-based transport overlay for management traffic with a distinct, application-layer policy enforcement point that performs granular command authorization. MTD and channel security are treated as monolithic solutions, not as complementary layers in a defense-in-depth strategy.
% \end{enumerate}


\section{Chapter Summary}

This chapter reviewed the state-of-the-art in network orchestration and security for disaggregated networks. We observed that while orchestration platforms are mature, they lack secure, granular integration with the data plane. Furthermore, existing defenses secure the communication channel but fail to enforce the specific command-level authorization required by Zero Trust architectures. Consequently, a critical gap exists for a unified framework that combines secure transport with fine-grained policy enforcement. The following chapter details the design of such a framework, which constitutes the core contribution of this thesis.




\newpage
\chapter{Proposed Method: Secure Framework Design and Architecture}
\label{chap:method}

Managing disaggregated network infrastructure, especially when distributed across multiple physical locations and cloud environments (e.g., AWS, on-premise data centers), presents a profound security challenge. Securing the data plane is not enough; the management plane itself, which carries sensitive configuration commands, becomes a high-value attack vector. A compromised management channel could allow an attacker to take control of the entire network fabric.

To address this, this thesis proposes a novel, layered, defense-in-depth architecture that implements Zero Trust principles at two distinct layers. The proposed method combines a \textbf{Zero Trust Overlay Network (OpenZiti)} for secure, identity-based transport with the \textbf{Secure Framework} developed in this work, which acts as a granular Zero Trust Policy Enforcement Point at the application layer. This multi-layered security model ensures that management traffic is protected both in transit across untrusted networks and at its destination before any action is taken.

\section{System Architecture Overview}
\label{sec:approach}

The complete system architecture, illustrated in Figure~\ref{fig:system_architecture}, demonstrates the synergy between the OpenZiti overlay and the Secure Framework. This design provides a holistic solution for secure, multi-cloud network orchestration.

\begin{figure}[t]
    \centering
    \includegraphics[width=\textwidth]{system_architecture.png} % Ensure you upload the image and name it correctly
    \caption{High-level System Architecture Combining OpenZiti and the Secure Framework.}
    \label{fig:system_architecture}
\end{figure}

\paragraph{Orchestration Plane - ONAP}
At the top, the ONAP environment contains the \textbf{SDN-C (OpenDaylight)} controller, which is the source of management intent. A user or automated process initiates a configuration change, which is triggered within SDN-C. The user's own access to the ONAP platform is secured via a secure tunnel, ensuring that the initial trigger is from an authenticated and authorized operator.

\paragraph{Secure Management Transport Plane - OpenZiti}
The critical challenge of connecting the central SDN-C to distributed management endpoints (our Secure Frameworks) across untrusted networks like the internet is solved by the \textbf{OpenZiti Overlay Network}. Instead of exposing NETCONF ports to the internet, the SDN-C's management traffic is tunneled through the OpenZiti fabric.
\begin{itemize}
    \item The \textbf{Ziti Controller} acts as the trust broker for the overlay, authenticating all participants (the SDN-C and the remote sites hosting the Secure Frameworks) and authorizing the creation of secure communication channels.
    \item \textbf{Edge Routers} serve as the ingress and egress points for the overlay. The SDN-C's NETCONF traffic enters the overlay via a local Edge Router and is securely transported through the fabric to the Edge Router at the destination site (e.g., AWS).
\item \textbf{Fabric Routers} participate solely in the Ziti fabric mesh and are dedicated to forwarding tunneled traffic. They do not run an Edge Listener, do not intercept traffic from LAN/WAN interfaces, and do not emit service traffic. Their role is strictly to relay overlay traffic between routers and maintain the mesh. 
\end{itemize}

This layer effectively creates a \textit{dark} network for management traffic, making the management endpoints invisible and inaccessible from the public internet. It enforces the first layer of Zero Trust: only authenticated and authorized entities can even send traffic to the management endpoints.

\paragraph{Zero Trust Enforcement Plane - Secure Framework}
While OpenZiti secures the transport, the \textbf{Secure Framework} provides the second, application-specific layer of Zero Trust enforcement. Deployed at each remote location (AWS, internal network), it receives the NETCONF traffic from the OpenZiti overlay. Its role is to:
\begin{itemize}
    \item Terminate the mTLS session from the SDN-C, verifying the controller's cryptographic identity.
    \item Perform a granular RBAC check to authorize the specific operation requested in the NETCONF RPC.
    \item If authorized, translate the NETCONF command into a gNMI command.
    \item Send the final, authorized gNMI command to the local SONiC switch.
\end{itemize}
This component acts as the final PEP, ensuring that even if the transport is secure, the controller can only perform explicitly permitted actions.

\paragraph{Managed Data Plane - SONiC}
At the bottom, the \textbf{SONiC} switches represent the distributed data plane fabric. Their management interfaces (gNMI) are only exposed to the local Secure Framework, completely isolating them from any direct network access.

\paragraph{End-to-End Workflow}
In this integrated architecture, a configuration command flows from the user to the SONiC switch through a multi-stage Zero Trust pipeline. The OpenZiti overlay provides a secure, authenticated path across the internet, while the Secure Framework provides the final, granular authorization and translation at the edge. This separation of concerns creates a highly resilient and secure architecture for managing geographically distributed, disaggregated networks.

\section{Core Translation Logic}

\label{sec:translation_logic}

At the heart of the Secure Framework lies the Core Translation Logic, a sophisticated bidirectional translation engine responsible for bridging the semantic and syntactic gap between the NETCONF and gNMI protocol ecosystems. This component serves as the central intelligence of the adapter, ensuring that configuration intent expressed by the ONAP controller through NETCONF is accurately transformed and applied to SONiC devices via gNMI. Equally important, responses and configuration data retrieved from the network devices are translated back into NETCONF-compliant formats that the ONAP SDN-C controller expects.

The translation process encompasses more than simple format conversion. It requires understanding of YANG data model structures, proper handling of XML namespaces, normalization of protocol-specific values, and preservation of operational semantics across the protocol boundary. The engine is designed with extensibility in mind, allowing support for additional YANG models beyond the currently implemented OpenConfig ACL modules.

\subsection{YANG Model and Path Translation}

\label{sec:yang_mapping}

Although both NETCONF and gNMI utilize YANG as their foundational data modeling language, they employ fundamentally different data representations and addressing mechanisms. NETCONF relies on XML encoding with XPath-based path expressions and XML namespace declarations to identify and structure configuration elements. In contrast, gNMI uses a structured path format composed of path elements with optional key specifications, combined with JSON-IETF encoding for data payloads. The translation engine bridges these representations through intelligent parsing and reconstruction.

\paragraph{Namespace Handling}

A significant challenge in NETCONF-to-gNMI translation is the variability of XML namespace usage across different NETCONF clients. The same YANG module may be referenced with different namespace prefixes depending on the client implementation. The translation engine addresses this through namespace-agnostic parsing techniques that identify elements by their local names rather than fully qualified names. This approach ensures robust processing regardless of the specific prefix conventions employed by the ONAP controller or other NETCONF clients.


\paragraph{Path Reconstruction}

The translation engine transforms XPath expressions into gNMI-compatible path formats while preserving the hierarchical relationships and list key specifications defined in the YANG models. For configuration elements that are list entries, such as ACL sets identified by name and type or ACL rules identified by sequence number, the engine extracts key values from the XML structure and formats them according to gNMI path conventions.

This path translation is model-aware, meaning the engine understands the structure of the OpenConfig ACL YANG model and can correctly identify which elements serve as list keys. This semantic understanding ensures that the translation produces valid gNMI paths that the SONiC device can correctly interpret and process.

\subsection{Protocol Operation Translation}

\label{sec:protocol_translation}

Beyond structural translation, the engine must map the operational semantics of NETCONF RPC to their gNMI equivalents. Each NETCONF operation carries specific meaning regarding how configuration should be retrieved or modified, and this meaning must be accurately conveyed through the corresponding gNMI operations. The Secure Framework implements translation for the primary configuration operations essential for network policy management, as detailed in Table~\ref{tab:protocol_mapping}.

\begin{table}[H]
    \centering
    \caption{Mapping of NETCONF RPCs to gNMI Operations}
    \label{tab:protocol_mapping}
    \begin{tabular}{|l|l|l|}
        \hline
        \textbf{NETCONF Operation} & \textbf{gNMI Operation} & \textbf{Description} \\
        \hline
        \texttt{<get-config>} & \texttt{Get} & Retrieves configuration data from a device. \\
        \texttt{<get>} & \texttt{Get} & Retrieves operational and configuration data. \\
        \texttt{<edit-config>} (merge) & \texttt{Set} (update) & Creates or modifies configuration. \\
        \texttt{<edit-config>} (delete) & \texttt{Set} (delete) & Removes configuration elements. \\
        \hline
    \end{tabular}
\end{table}

\paragraph{Configuration Retrieval}

When the framework receives a NETCONF \texttt{<get-config>} request from the ONAP controller, the translation engine initiates a corresponding gNMI Get operation targeting the relevant configuration paths on the SONiC device. The gNMI response, containing configuration data encoded in JSON-IETF format according to OpenConfig YANG model structures, undergoes reverse translation into XML.

This reverse translation iterates through the JSON response hierarchy, constructing equivalent XML elements with proper namespace declarations. For ACL configurations, this includes translating ACL set metadata such as names, types, and descriptions, as well as individual rule parameters including match criteria and forwarding actions. The resulting XML structure is wrapped in the standard NETCONF \texttt{<rpc-reply>} envelope, presenting the ONAP controller with a response indistinguishable from one that would be received from a native NETCONF device.

\paragraph{Configuration Modification}

The translation of \texttt{<edit-config>} operations requires careful handling of modification semantics. NETCONF supports multiple operation types including merge for adding or updating configuration, replace for substituting entire subtrees, and delete for removing configuration elements. The engine inspects operation attributes in the incoming XML to determine the intended modification type.

For create and merge operations, the engine extracts all relevant configuration parameters from the XML payload. In the context of ACL policy management, this encompasses table-level attributes such as name and type, rule-level parameters including sequence identifiers, protocol specifications, address and port match criteria, and forwarding actions. The engine also processes interface binding information that specifies policy attachment points and enforcement direction.

Value normalization forms a critical part of this extraction process. Protocol identifiers may appear in various formats across different systems, including numeric IP protocol numbers, OpenConfig enumeration values, or plain text representations. The translation engine maintains mapping tables to convert between these representations, ensuring that values are properly formatted for the target gNMI interface regardless of their original format in the NETCONF request.

For delete operations, indicated by the presence of a delete operation attribute in the source XML, the engine constructs gNMI Set requests containing delete paths. These paths precisely identify the configuration elements to be removed, whether entire ACL tables or specific rules within a table.

\paragraph{Transactional Integrity}

A key design principle of the translation engine is preservation of transactional semantics. When the ONAP controller issues a single \texttt{<edit-config>} request containing multiple interdependent configuration elements, such as an ACL table definition along with its associated rules and interface bindings, the engine consolidates these into a single gNMI Set request.

The gNMI protocol inherently supports atomic multi-path operations, allowing the SONiC device to apply all specified changes as a single transaction. This atomicity prevents partial application of configuration that could leave the device in an inconsistent state. If any portion of the configuration fails validation or application, the entire operation is rejected, and the device configuration remains unchanged.

Response translation maintains this transactional view. Successful gNMI operations result in standard NETCONF \texttt{<ok/>} responses, while failures are mapped to appropriate NETCONF error structures containing error type classifications, standardized error tags, and descriptive messages. This bidirectional translation of both requests and responses provides the ONAP controller with a consistent NETCONF interaction experience while leveraging the efficiency and capabilities of gNMI for actual device communication.



\section{Integrated Security Module}
\label{sec:security_module}
A fundamental differentiator of the proposed framework from existing protocol gateways is its integrated security module. This module is not an afterthought but a core component of the request processing pipeline, designed to enforce the principles of Zero Trust at every stage of a management operation. It transforms the framework from a simple translator into an active PEP. The complete security and translation workflow is detailed in Figure~\ref{fig:detailed_workflow}.

\begin{figure}[H]
    \centering
    \includegraphics[width=0.8\textwidth]{detailed_workflow-2.png} % 
    \caption{Detailed Internal Workflow of the Secure Framework.}
    \label{fig:detailed_workflow}
\end{figure}




\subsection{Secure Channel Establishment}
\label{sec:secure_channel}

The proposed architecture implements a mTLS framework to ensure end-to-end encryption across heterogeneous network management protocols. As illustrated in Figure~\ref{fig:detailed_workflow}, secure communication is established through two distinct protocol boundaries: the northbound NETCONF interface and the southbound gNMI interface.

The NETCONF server operates on port 6513, accepting secure connections from the SDN Controller (SDN-C) using the standard \texttt{<rpc>} and \texttt{<rpc-reply>} message exchange pattern. Each connection is authenticated through X.509 certificates issued by a unified Certificate Authority, ensuring mutual authentication between the controller and the adapter layer. The certificate hierarchy follows the PKIX standard, with role and identity information bound to the certificate (e.g., Common Name, Subject Alternative Name or Organizational Unit), enabling identity-based access decisions at the protocol level.”

On the southbound interface, the gNMI client establishes gRPC connections to the target network element on port 50051. This connection utilizes separate client certificates, allowing for credential isolation between the management layer and the device layer. The dual-certificate approach provides cryptographic separation of concerns, ensuring that compromise of northbound credentials does not directly expose southbound device access.

\subsection{Granular Access Control with Service Account Manager}
\label{sec:granular_access}

The Service Account Manager serves as the central authorization component within the adapter architecture, implementing RBAC for fine-grained permission management. As depicted in the architectural flow (steps 2-4 in Figure~\ref{fig:detailed_workflow}), each incoming NETCONF request triggers an authorization sequence before any configuration operation is executed.

Upon receiving a configuration request, the NETCONF server extracts the authenticated identity from the TLS session and initiates an access verification request to the Service Account Manager (step 2). The manager evaluates the request against the configured policy mappings, which define the permissible role transitions between the ONAP domain and the SONiC domain. Two principal roles are defined within this framework:

\begin{itemize}
    \item \textbf{Administrator}: Possesses full read-write access to all device configuration paths, including security-sensitive parameters and system-level settings.
    \item \textbf{Operator}: Restricted to read operations and a limited subset of write operations, typically encompassing operational state queries and non-critical configuration modifications.
\end{itemize}

The policy engine supports dynamic role assumption, wherein an authenticated administrator may optionally operate with reduced privileges when interacting with specific device contexts. Upon successful authorization (step 3), the Service Account Manager issues a role grant (step 4), which instructs the gNMI client to utilize the appropriate credential set for the downstream connection. This assume-role mechanism provides operational flexibility while maintaining the principle of least privilege.

\subsection{State Management and Translation Execution}
\label{sec:state_management}

The Adapter Core component orchestrates the semantic translation between YANG-modeled NETCONF operations and the Protocol Buffers-based gNMI interface. As shown in step 4 of Figure~\ref{fig:detailed_workflow}, the NETCONF server delegates parsed XML payloads along with their corresponding YANG schema context to the Adapter Core for processing.

The translation pipeline operates in three phases. First, the incoming NETCONF \texttt{<edit-config>} payload is validated against the applicable YANG modules to ensure structural conformance and constraint satisfaction. Second, the validated data tree is transformed into the equivalent gNMI path-value representation, mapping YANG container hierarchies to gNMI path elements and leaf values to typed scalar encodings. Third, the translated request is serialized into Protocol Buffer format (step 5) and dispatched to the gNMI client for transmission.

The gNMI client maintains persistent connections to the target SONiC device, utilizing the Translib framework on the device side for configuration application. Translib serves as the southbound abstraction layer, translating gNMI Set and Get operations into the underlying Redis database transactions that constitute the SONiC configuration store. This architectural separation enables the adapter to remain agnostic to device-specific implementation details while leveraging the standardized gNMI interface for configuration management.

Response handling follows the inverse path, with gNMI response messages being translated back into NETCONF \texttt{<rpc-reply>} format for delivery to the originating SDN-C (step 6). Error conditions encountered at any layer are mapped to appropriate NETCONF error elements, preserving semantic fidelity across the protocol boundary.

\section{Chapter Summary}
\label{sec:method_summary}

This chapter detailed a Secure Framework designed to enforce Zero Trust security in disaggregated networks by acting as a trusted intermediary between the ONAP orchestrator and the SONiC data plane. The architecture provides a high-fidelity translation layer that bridges the protocol gap between NETCONF and gNMI, while simultaneously implementing a multi-layered security pipeline grounded in the NIST ZTA model. It leverages an OpenZiti overlay network to create secure, isolated transport channels and enforces command-level RBAC to ensure least-privilege access for every management operation. In essence, it is presented not just as a protocol gateway, but as a comprehensive security solution that combines protocol mediation with a defense-in-depth architecture, providing the necessary foundation for automated and secure network management.





\newpage
\chapter{Experiments and Evaluation}
\label{chap:experiments}

This chapter presents a comprehensive evaluation of the proposed Zero Trust policy enforcement framework. The core of this thesis is not just a single component, but a multi-layered security architecture that combines an identity-based secure overlay for transport with a purpose-built PEP for command-level authorization. Consequently, a simple performance benchmark is insufficient to validate our contributions. Instead, this chapter introduces and applies a \textbf{hybrid evaluation methodology}, as outlined in our research challenges, to rigorously assess the framework's security posture and architectural novelty.


% \section{Experimental Environment}
% \label{sec:exp_environment}

\section{Experiment design overview}
\label{sec:eval_methodology}

As outlined in the research challenges, we evaluate a security architecture whose primary contributions are granular policy enforcement and architectural isolation. This approach combines:

\begin{enumerate}
    \item Perform an \textbf{Experimental Security Validation} to provide empirical evidence that the implemented security controls function correctly and effectively enforce the defined policies.
    
    \item Conduct a \textbf{Comparative Architectural Analysis} to demonstrate the framework's design superiority in mitigating a wide range of threats compared to existing single-paradigm security solutions.

\end{enumerate}
This dual-pronged evaluation directly addresses the core research questions regarding the framework's architectural completeness, its ability to enforce policy, and its functional correctness in a practical implementation. To assess the architectural completeness of our solution, we use the \textbf{STRIDE threat model} ~\cite{Khan2017stride, ruffy2016stride}. STRIDE, an acronym for Spoofing, Tampering, Repudiation, Information Disclosure, Denial of Service, and Elevation of Privilege, provides a structured framework for identifying and categorizing security threats. 

\section{Testbed and implementation details}

The experimental testbed was deployed on a virtualized cluster emulation environment running Ubuntu Server 20.04 LTS. ONAP master and worker nodes are interconnected via a private management network to ensure isolation. Table~\ref{tab:hardware_specs} details the specific resource allocation and roles for each component in the cluster.

\begin{table}[t]
\centering
\small
\caption{Testbed Hardware Configuration and Resource Allocation}
\label{tab:hardware_specs}
\begin{tabular}{p{2.4cm} c p{5cm} p{7.5cm}}
\toprule
\textbf{Node Role} & \textbf{Qty.} & \textbf{Resources (vCPU / RAM / SSD)} & \textbf{Function \& Responsibility} \\ \midrule
ONAP Master & 1 & 32 Cores / 96 GB / 200 GB & Hosts core orchestration services and control plane components. \\ \addlinespace
ONAP Worker & 2 & 8 Cores / 16 GB / 200 GB & Executes distributed microservices; ensures scalability and fault isolation. \\ \addlinespace
SONiC Node & 2 & 12-16 Cores / 16 GB / 200 GB & Simulates the data plane (3S-NOS) \\

\addlinespace
Ziti Controller & 1 & 2 Cores / 4 GB / 30 GB & Manages the Ziti router and secure zero-trust connectivity between system components. \\ \bottomrule
\end{tabular}
\end{table}


The orchestration plane operates on the \textbf{ONAP Paris} release, with the management layer deployed as containerized microservices within a Kubernetes cluster. The secure transport plane utilizes \textbf{OpenZiti version 1.6.9} to provide protected communication channels. The \textbf{Secure Framework} is implemented using \textbf{Python version 3.8}. The data plane consists of two \textbf{SONiC-VS} instances, we built images at version \textbf{2025-09}.



\section{Evaluation Metrics}
\label{sec:evaluation_metrics}
Given that our proposed method is a security architecture whose primary contributions are granular policy enforcement and architectural isolation, a purely quantitative performance benchmark is insufficient to validate its effectiveness. Instead, our evaluation focuses on the framework's ability to mitigate specific, well-defined security threats. To provide a systematic and comprehensive analysis, we adopt the \textbf{STRIDE threat model} as our evaluation framework~\cite{Khan2017stride, ruffy2016stride}.

STRIDE is a widely recognized methodology developed by Microsoft for identifying and categorizing security threats. The model's name is an acronym for the six fundamental threat categories it addresses:

\begin{itemize}
    \item \textbf{Spoofing:} Illegally assuming the identity of another user, device, or component.
    \item \textbf{Tampering:} Maliciously modifying data, either in transit or at rest.
    \item \textbf{Repudiation:} Denying having performed an action, where the system cannot prove otherwise.
    \item \textbf{Information Disclosure:} Exposing information to individuals not authorized to access it.
    \item \textbf{Denial of Service:} Preventing legitimate users from accessing a service or resource.
    \item \textbf{Elevation of Privilege:} Gaining capabilities beyond what is authorized for a given identity.
\end{itemize}

By using STRIDE, we can systematically analyze how architectural design and specific components address each of these threat categories. This approach allows us to move beyond simple performance metrics and instead evaluate the framework based on the security properties it is designed to guarantee. The subsequent experimental validation in Section \ref{sec:scenarios} will demonstrate efficacy in mitigating threats from each of these categories in practical scenarios.

\section{Scenarios and Results}
\label{sec:scenarios}
To evaluate the effectiveness of our solution in enforcing Zero-Trust policies in disaggregated network environments, we performed several experiments based on the STRIDE threat model ~\cite{Khan2017stride, ruffy2016stride}. Each experiment targeted a specific scenario, analyzed the potential attack surface, and observed the system's response to simulated threats. The main objectives were to assess identity-based access control, consistent policy enforcement across protocols, and the system's resilience during runtime.

\subsection{Scenario A: Client connection with Invalid/Revoked Identity}

\paragraph{Objective} To verify that identity-based overlay network prevents unauthorized entities from gaining access to the management plane's PEP.

\paragraph{Threat Procedure} An attacker, operating from the Untrusted Zone, attempts to connect directly to the Policy Engine's endpoint. The attack is simulated using two methods: (1) attempting to use a stolen but expired enrollment token for the OpenZiti overlay, and (2) replaying a previously valid enrollment request.

\paragraph{Expected Outcome} In both cases, the OpenZiti controller at the edge of the overlay network rejects the connection attempt during the initial handshake . The attacker's requests never reach the Policy Engine. 




\begin{figure}[t]
    \centering
    \includegraphics[width=0.5\linewidth]{openziti-authentication.png}
    \caption{OpenZiti overlay network rejecting unauthorized identity enrollment.}
    \label{fig:openziti-authentication}
\end{figure}



\subsection{Scenario B: Malicious command injection from authenticated source}

\paragraph{Objective} To demonstrate that the Policy Engine prevents the execution of maliciously modified management commands, even if they originate from a legitimately authenticated source.

\paragraph{Threat Procedure} This test simulates an attacker who has compromised a legitimate orchestrator (e.g., ONAP). The attacker intercepts an authorized NETCONF \texttt{<edit-config>} command destined for a switch and modifies its payload-for example, creating a firewall rule  'allow/deny'. The tampered command is then sent through the established secure tunnel.

\paragraph{Expected Outcome} The Policy Engine receives the command from the authenticated (but compromised) orchestrator. It validates the command's content against the granular policy associated with the orchestrator's role. The engine detects that the requested action (e.g., creating an allow/deny rule for a sensitive port) is not permitted for that entity or context. The command is rejected, and a high-severity alert is generated in the centralized logger, flagging the tampering attempt and the source identity.



\begin{figure}[H]
    \centering
    \includegraphics[width=\linewidth]{detection-logs.png}
    \caption{Secure Framework Detection Logs}
    \label{fig:detection_logs}

\end{figure}



\subsection{Scenario C: Valid Identity attempting cross-segment access}
\label{sec:scenario-c}

\paragraph{Objective} To validate that the identity-based overlay network enforces security segmentation at the transport layer, preventing a legitimately authenticated entity from accessing management endpoints outside of its explicitly authorized security segment.

\paragraph{Threat Procedure} This scenario simulates an orchestrator 3S-COM (ONAP) attempting to manage devices across different security segments. The 3S-COM identity was created and granted access only to the AWS service segment by OpenZiti Edge Router Policies, shown in the bottom half of \autoref{fig:onap-ziti-policy}. An 3S-COM entity using this authenticated identity then attempted to initiate a NETCONF request targeting a router whose management endpoint is part of the internal service segment. 

\paragraph{Expected Outcome}
The connection attempt from the ONAP orchestrator succeeds to 3S-NOS (SONiC) on AWS. However, the connection attempt to 3S-NOS on internal is rejected at the transport layer by the OpenZiti overlay. The management request never reaches internal switch itself. The ONAP UI accurately reflects this outcome, showing \textbf{sonic-aws} as "Connected" and \textbf{sonic-internal} as "UnableToConnect," thus empirically validating the successful enforcement of transport-layer segmentation.

\begin{figure}[H]
    \centering
    \includegraphics[width=\columnwidth]{onap-ziti-policy.png}
    \caption{Validation of transport-layer segmentation.}
    \label{fig:onap-ziti-policy}
\end{figure}




\subsection{Scenario D: Unauthorized operation from read-only role}

\paragraph{Objective:}
To verify that command-level policy engine prevents an entity from escalating its privileges by executing write operations when it is only authorized for read-only access.

\paragraph{Threat Procedure:}
 An attacker compromises an operator account that has legitimate read-only access for monitoring purposes. The attacker then attempts to leverage this authenticated session to perform a write operation, effectively trying to elevate their privileges from passive observer to active administrator.

\paragraph{Procedure:}
A user was authenticated with a \textit{read-only} role, which permits \texttt{<get>} and \texttt{<get-config>} operations but explicitly denies \texttt{<edit-config>}. The authenticated user sent a syntactically correct but unauthorized NETCONF \texttt{<edit-config>} command to modify a device's ACL  configuration. We observed the Policy Engine's evaluation of the specific RPC operation against the user's role and the resulting action.




\begin{figure}[H]
    \centering
    \includegraphics[width=\linewidth]{access-write.png}
    \caption{Denial of an unauthorized write operation}
    \label{fig:access-deny-write}
\end{figure}


\subsection{Scenario E: Valid policy deployment by authorized administrator}
\label{sec:scenario-e}

\paragraph{Objective}
To validate the end-to-end policy pipeline, from intent translation to data plane enforcement, by deploying a specific ACL rule and verifying its successful installation on the target SONiC switch.

\paragraph{Threat Procedure}
This scenario represents a standard, authorized operational workflow. An administrator with the appropriate privileges defines a policy to allow TCP traffic on port 80 from a specific source IP range to a destination IP range. The policy is submitted to the framework for translation and enforcement.

\paragraph{Expected Outcome}
The Policy Engine authenticates the administrator and authorizes the request based on their role. The Secure Framework converts the abstract policy into a concrete gNMI `Set` command for the SONiC device. The command is successfully executed, and the new ACL rule is installed in the data plane. The framework's user interface reports a successful deployment, and a direct verification on the switch's command-line interface confirms that the rule is active and correctly configured. This demonstrates the framework's ability to reliably translate and enforce authorized policies.


\begin{figure}[H]
    \centering
    \includegraphics[width=\columnwidth]{deployment_result.png}
    \caption{Successful deployment of an ACL rule.}
    \label{fig:deployment-result}
\end{figure}


The collective results from these four scenarios demonstrate that our methods successfully mitigates the targeted STRIDE threats through its multi-layered defense strategy. To provide a comprehensive overview of the evaluation, \autoref{tab:test_case_eval} summarizes the functional validation results, explicitly mapping each experimental scenario to its specific test case, targeted threat category, and verified outcome.

\begin{table}[H]
\centering
\caption{Functional Validation Results Security Controls}
\label{tab:test_case_eval}
\begin{tabular}{c p{6cm} c c c}
\toprule
\textbf{Ex.} & \textbf{Specific Test Case / Attack Simulation} & \textbf{STRIDE Ref.} & \textbf{Exp.} & \textbf{Out.} \\ 
\midrule
A & Client connection with Invalid/Revoked Identity & Spoofing & Block & $\checkmark$ \\ 
\addlinespace
B & Malicious command injection from authenticated source & Tampering, Repudiation & Deny & $\checkmark$ \\ 
\addlinespace
C & Valid Identity attempting cross-segment access & Info. Disclosure & Block & $\checkmark$ \\ 
\addlinespace
D & Unauthorized operation from read-only role & EoP & Deny & $\checkmark$ \\ 
E & Valid policy deployment by authorized administrator & – & Allow & $\checkmark$ \\ 
\bottomrule
\multicolumn{5}{l}{\footnotesize \textit{Exp. = Expected Action; Out. = Actual Outcome; EoP = Elevation of Privilege.}}
\end{tabular}
\end{table}






\section{Evaluation}
\label{sec:resultEvaluation}
\subsection{Results Analysis}
\label{sec:results_analysis}
The experimental results presented in Section \ref{sec:scenarios} provide empirical validation for the architecture and directly address the key research questions. The analysis confirms that the framework's multi-layered design effectively enforces Zero-Trust principles in a heterogeneous, disaggregated environment.

Regarding \textbf{RQ1}, which questions how to enforce fine-grained policies in dynamic environments, our evaluation demonstrates a robust solution. Scenarios B and D validate that command-level Policy Engine moves beyond simple entity authentication. By inspecting every management command, the framework successfully prevented a compromised but authenticated orchestrator from executing a malicious write operation (Scenario B) and blocked an authenticated read-only user from attempting to elevate their privileges (Scenario D). This confirms that our proposed method enforces the principle of least privilege at the operational level, ensuring that even authenticated entities cannot misuse their access.

The challenge of \textbf{cross-protocol policy consistency (RQ2)} is addressed by mediating architecture. The successful end-to-end policy deployment in Scenario E, which required translating an abstract policy into a concrete gNMI command for the SONiC data plane, empirically validates this capability. By channeling all management requests through the PEP, the framework ensures that a single, unified security policy is applied, regardless of whether the originating protocol is NETCONF or gNMI. This design eliminates the risk of policy gaps or inconsistencies that could arise from separate, uncoordinated protocol handlers.

\textbf{RQ3} focuses on achieving architectural isolation for management endpoints. Our results from Scenarios A and C demonstrate the efficacy of identity-based overlay network. In Scenario A, an external attacker without a valid cryptographic identity was blocked at the transport layer, unable to even discover the Policy Engine's endpoint. More powerfully, Scenario C showed that even a legitimately authenticated entity was prevented from accessing a management endpoint outside its authorized security segment, with the connection status correctly reflected in the ONAP UI. This confirms that the framework successfully implements an \textit{authenticate before connect} workflow, effectively isolating management interfaces and preventing unauthorized reconnaissance and access.

Finally, this work addresses \textbf{RQ4} by presenting a methodology for evaluating a multi-layered Zero-Trust architecture. Instead of relying on simple performance metrics, we combined a qualitative STRIDE-based threat analysis with empirical, scenario-based validation on a functional prototype. This hybrid approach allowed us to systematically map architectural components to specific threat mitigations (as shown in \autoref{tab:stride_comparison}) and then provide concrete evidence that these mitigations function as designed in practical use cases (Scenarios A-E). This methodology provides a comprehensive and verifiable assessment of the framework's security posture.

\subsection{Comparison to State-of-the-Art Methods}
\label{sec:comparisonSOTA}




We apply the STRIDE framework to qualitatively assess its ability to mitigate threats specific to the management plane of disaggregated networks. \autoref{tab:stride_comparison} evaluates proposed framework against related works, demonstrating that while existing solutions address certain threats within their specific domains, our proposed method provides a more comprehensive defense by focusing on the management plane as a unified attack surface. Its mediating architecture and identity-based overlay are designed to counter threats that exploit protocol inconsistencies and implicit trust.

The analysis in \autoref{tab:stride_comparison} highlights prevention-first approach to management plane security. Its strength lies in the architectural integration of an identity-based overlay network with a mediating PEP. This design fully prevents \textbf{Spoofing} by ensuring every management entity is cryptographically authenticated before it can connect, a robust approach shared by other identity-centric solutions like TaaS~\cite{Erel-Ozcevik2024}, IBC for SDN~\cite{lam2015securing}, and Bloccess~\cite{Ding2022}.

Against \textbf{Tampering}, our framework provides strong prevention by validating every command against authorized policies before execution. This proactive interception is a key differentiator from blockchain-based frameworks like B-DAC~\cite{duy2022} and Bloccess~\cite{Ding2022}, which offer tamper-evidence through an immutable ledger but do not prevent the initial malicious command. Similarly, while IBC for SDN~\cite{lam2015securing} ensures in-transit integrity, it is blind to the content of the commands themselves, a gap that our framework command-level PEP explicitly closes.
\clearpage
\begin{table}[H]
    \centering
    \caption{STRIDE-Based Threat Mitigation Efficacy Comparison}
    \scriptsize
    \label{tab:stride_comparison}
    \begin{tabular}{@{}p{2cm}l p{1.2cm}c p{1cm}c p{1cm}c p{1cm}c p{1cm}c p{1cm}c p{1cm}c p{1cm}c@{}}
        \toprule
        \textbf{STRIDE Threat} & \textbf{Ours} & \textbf{TaaS~\cite{Erel-Ozcevik2024}} & \textbf{6G RANZT~\cite{Sedjelmaci2024}} & \textbf{IBC for SDN~\cite{lam2015securing}} & \textbf{P4 DDoS~\cite{Musumeci2021}} & \textbf{B-DAC~\cite{duy2022}} & \textbf{MTD~\cite{mendonca2022security}} & \textbf{Bloccess~\cite{Ding2022}} \\
        \midrule
        \textbf{Spoofing} & \checkmark\checkmark\checkmark\checkmark & \checkmark\checkmark\checkmark\checkmark & \checkmark\checkmark\checkmark & \checkmark\checkmark\checkmark\checkmark & N & \checkmark\checkmark\checkmark & \checkmark & \checkmark\checkmark\checkmark\checkmark \\
        \addlinespace
        \textbf{Tampering} & \checkmark\checkmark\checkmark\checkmark & \checkmark\checkmark\checkmark & \checkmark\checkmark & \checkmark\checkmark\checkmark & N & \checkmark\checkmark\checkmark\checkmark & N & \checkmark\checkmark\checkmark\checkmark \\
        \addlinespace
        \textbf{Repudiation} & \checkmark\checkmark\checkmark & \checkmark & \checkmark & N & N & \checkmark\checkmark\checkmark\checkmark & N & \checkmark\checkmark\checkmark\checkmark \\
        \addlinespace
        \textbf{Information Disclosure} & \checkmark\checkmark\checkmark\checkmark & \checkmark\checkmark\checkmark & \checkmark\checkmark\checkmark & \checkmark\checkmark\checkmark\checkmark & N & \checkmark\checkmark & \checkmark\checkmark & \checkmark\checkmark\checkmark \\
        \addlinespace
        \textbf{Denial of Service} & \checkmark\checkmark & \checkmark\checkmark & \checkmark\checkmark & \checkmark & \checkmark\checkmark\checkmark & \checkmark\checkmark & \checkmark\checkmark & \checkmark\checkmark \\
        \addlinespace
        \textbf{Elevation of Privilege} & \checkmark\checkmark\checkmark\checkmark & \checkmark\checkmark\checkmark & \checkmark\checkmark\checkmark & \checkmark\checkmark\checkmark & N & \checkmark\checkmark & N & \checkmark\checkmark\checkmark\checkmark \\
        \bottomrule
        \multicolumn{9}{l}{\scriptsize{\checkmark\checkmark\checkmark\checkmark: Robust Prevention \quad \checkmark\checkmark\checkmark: Mitigation \quad \checkmark\checkmark: Partial Mitigation \quad \checkmark: Limited Mitigation \quad N: Not Addressed}}
    \end{tabular}
\end{table}

% \subsubsection{Threat Mitigation Efficacy: STRIDE}






For \textbf{Information Disclosure}, our a offers a superior defense by making management endpoints \textit{dark} via its overlay network, a preventative measure that Moving Target Defense (MTD) approaches~\cite{mendonca2022security} only mitigate through obfuscation. Against \textbf{Denial of Service}, our dedicated telemetry and automated response components provide a robust \textbf{Detect \& Respond} capability for the management plane, complementing data-plane-focused solutions like P4-based DDoS detection~\cite{Musumeci2021} which do not protect control and management interfaces.

Finally, our proposed method is uniquely effective against \textbf{Elevation of Privilege}. Unlike session-based models where an authenticated entity gains broad permissions, command-level RBAC enforces least privilege on every single operation. This prevents an authenticated but compromised orchestrator from performing actions beyond its explicit mandate-a critical vulnerability not fully addressed by most other compared methods, with the exception of highly secure, blockchain-enforced systems like Bloccess~\cite{Ding2022}.


% \section{Experimental Security Validation}
% \label{sec:experimental_validation}
% \subsection{Experimental Results}
% \label{sec:results}

\section{Chapter Summary}
\label{sec:related_summary}
This chapter presented a robust, hybrid evaluation of the proposed Secure Framework. Through a comparative architectural analysis against the STRIDE threat model, we established the framework's design superiority over existing security mechanisms by demonstrating its ability to provide defense-in-depth. We then provided empirical proof of the system's correctness through a series of security validation experiments. These tests confirmed that the framework successfully enforces Zero Trust principles by combining transport-layer isolation with application-layer, command-level policy enforcement. The evaluation concludes that the proposed framework is a novel, effective, and practical solution for enforcing Zero Trust policies in complex, disaggregated network environments.

\chap{Conclusion and Future Works}
\label{chap:conclusions}

This thesis successfully addressed the critical security challenges in disaggregated networks by designing, implementing, and evaluating the Secure Framework, a purpose-built PEP for Zero Trust policy enforcement. The framework resolves the four research questions through: (1) fine-grained, command-level RBAC enforcement ensuring least privilege access; (2) protocol-agnostic translation between NETCONF and gNMI maintaining policy consistency; (3) identity-based transport security architecturally isolating management endpoints; and (4) a hybrid evaluation methodology combining STRIDE threat modeling with empirical validation. The results demonstrate that integrating defense-in-depth security controls with granular policy enforcement provides a functionally correct, practical solution for implementing never trust, always verify principles in automated network environments.

\sect{Main Contributions}
\label{sec:contributions}

This thesis makes several key contributions to the field of network security and automation:

\begin{enumerate}
    \item \textbf{A Novel Policy Enforcement Framework for Zero Trust:} We presented the complete design and architecture of the \textbf{Secure Framework}. Its primary innovation is the integration of a RBAC engine that performs fine-grained, command-level authorization for every management request. This enforces the principle of least privilege on the orchestration system itself, a critical capability missing from existing solutions.

    \item \textbf{An Architectural Pattern for Securing Disaggregated Networks:} This work demonstrated a practical, multi-layered architectural pattern for implementing ZTA. By integrating the Secure Framework with an identity-based overlay network (OpenZiti), we showed how to achieve defense-in-depth. This pattern solves two distinct but related problems: it provides transport-layer security through architectural isolation (making management endpoints ``dark'') and application-layer security through granular policy enforcement.

    \item \textbf{A Hybrid Evaluation Methodology for Security Architectures:} We developed and applied a novel, hybrid evaluation methodology tailored for assessing policy-centric security frameworks. This approach combines a qualitative \textbf{Comparative Architectural Analysis} using the STRIDE threat model with a quantitative \textbf{Experimental Security Validation}. This methodology provides a rigorous means to demonstrate both the design superiority and the functional correctness of a security architecture, a contribution to the practice of security research itself.
\end{enumerate}
The results in this thesis have also been used to write the article \textit{ZEP-DN: A Zero-Trust Policy Enforcement Framework for Automated Management in Disaggregated Network Systems}, which has been submitted for publication to the \textbf{Journal of Network and Systems Management} [TVS1].

\sect{Limitations}
\label{sec:limitations}

While this thesis successfully demonstrates a robust and novel security architecture, the scope of the work has several limitations that present opportunities for future research:

\begin{enumerate}
    \item \textbf{Focus on Security over Performance:} The evaluation was intentionally focused on validating the architectural design and the functional correctness of the security controls. A comprehensive quantitative performance analysis - measuring end-to-end latency, transaction throughput under load, and scalability to thousands of devices—was beyond the scope of this work.
    
    \item \textbf{Specific Protocol and Use Case:} The framework was designed and implemented for the specific \textbf{NETCONF $\rightarrow$ gNMI} translation path required for the ONAP-SONiC integration use case. Support for other protocols (e.g., RESTCONF) or the reverse translation direction was not addressed.
    
    \item \textbf{Static Policy Model:} The current implementation relies on a RBAC model. While effective, RBAC policies are relatively static. The framework does not yet incorporate more dynamic, context-aware authorization models like Attribute-Based Access Control.
    
    \item \textbf{Virtualized Evaluation Environment:} The experimental validation was conducted in a controlled, virtualized environment using SONiC-VS. The operational challenges of deploying and managing such a framework in a large-scale, physical data center, including at-scale certificate lifecycle management and integration with enterprise PKI systems, were not explored.
\end{enumerate}

\newpage
\sect{Future Works}
\label{sec:future_works}

The limitations of this work provide a clear and promising roadmap for future research. The following directions represent logical next steps to build upon the foundation established in this thesis:

\begin{enumerate}
    \item \textbf{Performance Benchmarking and Optimization:} A critical next step is to conduct a rigorous performance evaluation of the Secure Framework. This would involve measuring the latency overhead of the security pipeline and benchmarking its maximum throughput. The results would guide optimization efforts, such as caching policy decisions or implementing asynchronous processing, to ensure the framework is suitable for carrier-grade deployments.

    \item \textbf{Advanced Policy Enforcement with ABAC:} A significant enhancement would be to evolve the RBAC engine into a full-fledged Attribute-Based Access Control system. This would allow the framework to make more dynamic authorization decisions based on real-time context, such as the time of day, the security posture of the controller, or threat intelligence feeds, more closely aligning with the dynamic "Trust Algorithm" concept from the NIST ZTA model.

    \item \textbf{Enhanced Security Analytics and Closed-Loop Automation:} The Secure Framework is a natural chokepoint for all management traffic, making it an ideal sensor for security analytics. Future work could focus on exporting detailed logs of all authorized and, more importantly, denied requests to a Security Information and Event Management (SIEM) system. This could enable a closed-loop automation system where the framework could dynamically update its own policies in response to detected threats.

    \item \textbf{Expanding Protocol Support and Bidirectionality:} To broaden its applicability, the framework could be extended to support other management protocols. Furthermore, implementing the reverse \textbf{gNMI $\rightarrow$ NETCONF} translation path would enable use cases where modern, gNMI-native controllers need to manage legacy network devices, making the framework a universal, secure adapter for heterogeneous environments.
\end{enumerate}



\chap{Publication}

\noindent [TVS1] Son V. Tran, Hung K. Luu, Hoa N. Nguyen, \textit{ZEP-DN: A Zero-Trust Policy Enforcement Framework for Automated Management in Disaggregated Network Systems}, submitted to Springer Journal of Network and Systems Management, December 2025. 
    

\newpage
\addcontentsline{toc}{chapter}{Bibliography} 
% \bibliographystyle{plainnat}
\bibliographystyle{ieeetr}
\bibliography{citation}



\end{document}